\documentclass[11pt]{article}
\pdfobjcompresslevel=0

\usepackage[review]{acl}

\usepackage{amsmath}
\usepackage{amssymb}
\usepackage{mathtools}
\usepackage{bm}
\usepackage{amsthm}

\usepackage{booktabs}
\usepackage{multirow}
\usepackage{array}
\usepackage{graphicx}
\usepackage{algorithm}
\usepackage{algpseudocode}
\usepackage{microtype}
\usepackage{enumitem}
\usepackage{xcolor}
\usepackage[nameinlink,capitalize]{cleveref}

% =========================================================
% Theorem environments
% =========================================================
\newtheorem{proposition}{Proposition}
\newtheorem{theorem}{Theorem}
\newtheorem{lemma}{Lemma}
\newtheorem{remark}{Remark}
\newtheorem{definition}{Definition}
\newtheorem{corollary}{Corollary}

% =========================================================
% CBG commands
% =========================================================
\newcommand{\R}{\mathbb{R}}
\newcommand{\Gh}{\mathcal{G}_{h}}
\newcommand{\Gb}{\mathcal{G}_{b}}
\newcommand{\Ghh}{\widehat{\mathcal{G}}_{h}}
\newcommand{\Gbh}{\widehat{\mathcal{G}}_{b}}
\newcommand{\Resc}{\rho_{\mathrm{esc}}}
\newcommand{\Barrier}{\mathcal{B}}
\newcommand{\U}{\mathcal{U}}
\newcommand{\Laln}{\mathcal{L}_{\mathrm{aln}}}
\newcommand{\Lbgn}{\mathcal{L}_{\mathrm{bgn}}}
\newcommand{\Lhfa}{\mathcal{L}_{\mathrm{hfa}}}
\newcommand{\rhohat}{\widehat{\rho}_{\mu}}
\newcommand{\rhoreg}{\rho_{\mu}}
\newcommand{\Lam}{\Lambda_{\theta}}
\newcommand{\Bmu}{\widehat{B}_{\mu}}
\newcommand{\Ft}{\widetilde{F}}
\newcommand{\Ch}{C_{h}}
\newcommand{\Cb}{C_{b}}
\newcommand{\etak}{\eta_{k}}
\newcommand{\Nb}{\mathcal{N}_{b}}
\newcommand{\Nh}{\mathcal{N}_{h}}
\newcommand{\diag}{\operatorname{diag}}
\newcommand{\row}{\operatorname{row}}
\newcommand{\tr}{\operatorname{tr}}
\newcommand{\Cth}{\widetilde{C}_{h}}
\newcommand{\Ctb}{\widetilde{C}_{b}}
\newcommand{\Pmu}{P_{\mu}}
\newcommand{\Smu}{S_{\mu}}

% Placeholder for results that do not exist yet. Never replace with a number
% that was not measured.
\newcommand{\tbd}{\textcolor{gray}{--}}
\newcommand{\todo}[1]{\textcolor{red}{[TODO: #1]}}


\title{\textbf{No Cheap Escape: Collateral Barrier Geometry for\\Self-Destructive Language Models}}
\author{Anonymous Authors}
\date{}

\begin{document}
\maketitle

\begin{abstract}
Safety alignment in open-weight language models can be removed by harmful fine-tuning or direct weight edits, often with little loss of capability. Self-destructive defenses make a particular harmful adaptation damage capability, but an adaptive attacker need not follow that direction and may find another that changes safety behavior at little benign cost. We propose the no-cheap-escape principle: a model should admit no admissible local parameter direction that causes a large safety-sensitive change at small benign cost. \textbf{Collateral Barrier Geometry (CBG)} instantiates this principle by comparing safety-sensitive and benign functional geometries and characterizing the weakest escape, the direction with the most safety change per unit of benign change, as a generalized eigenvalue problem. To make this trainable on minibatches, CBG damps the rank-deficient geometry with a rescaling-invariant metric and minimizes a normalized top-$k$ escape efficiency, computed exactly from a small Gram matrix of per-example gradients and differentiated at the cost of two additional gradient evaluations. The theory is local; we evaluate CBG empirically under standard, parameter-efficient and adaptive low-collateral fine-tuning attacks and under refusal-direction orthogonalization.
\end{abstract}



\section{Introduction}
\label{sec:introduction}

Safety alignment reduces harmful behavior in large language models (LLMs), but it often does not survive modification after release. Fine-tuning on a few harmful examples can remove refusal behavior while leaving general capability largely intact~\citep{qi2024finetuning}, and a direct weight edit can remove a single refusal direction~\citep{arditi2024refusal}. For open-weight models, alignment must therefore concern not only the released checkpoint but also how its behavior responds to subsequent parameter changes.

Many defenses try to make harmful fine-tuning harder~\citep{huang2024vaccine,huang2025booster,rosati2024repnoise,tamirisa2025tar}, yet several lose their protection under stronger attacks~\citep{qi2025evaluating}. Self-destructive alignment takes a complementary route~\citep{henderson2023self}: rather than only resisting harmful adaptation, it makes that adaptation damage benign capability. The closest prior work, SEAM~\citep{wang2026seam}, trains the model so that the gradient of a harmful loss opposes the gradient of a benign loss, so following the former degrades the latter.

Such coupling makes one harmful optimization direction destructive. It does not establish that every other way of changing safety behavior is also expensive. An adaptive attacker is not bound to the defender's harmful loss: it may change the objective, optimizer or trainable parameters, add a regularizer that preserves benign behavior, or edit weights directly. SEAM's own analysis leaves capability-preserving attacks as an open question and reports that refusal-direction orthogonalization remains effective against SEAM-defended models. This raises a sharper question: \emph{what is the cheapest way to change safety-sensitive behavior while preserving benign functionality?}

We make this question the defense target through the \emph{no-cheap-escape principle}: a self-destructive model should admit no admissible local parameter direction that produces a large safety-sensitive change at small benign cost. The principle shifts attention from one harmful trajectory to the direction most favorable to the attacker. If even that direction is expensive, no cheap local alternative exists; if it is not, protecting other directions matters little.

We instantiate the principle as \textbf{Collateral Barrier Geometry (CBG)}. For every admissible update direction, CBG compares how much the update moves the model's preference between safe and harmful responses with how much it moves behavior on benign inputs. The direction with the largest safety change per unit of benign change is the \emph{weakest escape}; its efficiency is a generalized eigenvalue of the two geometries and is inversely proportional to a \emph{collateral barrier}, the minimum modeled benign cost of a prescribed safety-sensitive change (\cref{fig:cbg-overview}). CBG trains the model to suppress its leading escapes. Doing so on minibatches requires a damped, rescaling-aware and normalized form of the efficiency, which we derive in \cref{sec:method} and which can be computed exactly and differentiated at the cost of two additional gradient evaluations per step.

\begin{figure*}[t]
    \centering
    \includegraphics[width=\textwidth]{CBG_method.pdf}
    \caption{CBG searches for low-collateral safety-changing directions and reshapes the local geometry so that leading escapes incur greater benign collateral.}
    \label{fig:cbg-overview}
\end{figure*}

The theory characterizes a local geometry and does not guarantee robustness to long fine-tuning trajectories or non-local weight edits; we test those settings empirically. Our contributions are:
\begin{itemize}[leftmargin=*]
\item We introduce the no-cheap-escape principle and formulate self-destructive alignment as a worst-direction benign-collateral problem, characterized spectrally by a collateral barrier.
\item We develop CBG, a trainable objective that remains well posed on minibatches, with an exact Gram-based spectral solver and an efficient gradient estimator.
\item We evaluate CBG under standard, parameter-efficient and adaptive low-collateral fine-tuning attacks and under orthogonalization, and analyze whether the learned geometry predicts attack outcomes.
\end{itemize}





\section{Related Work}
\label{sec:related}

\paragraph{Harmful fine-tuning and defenses.}
Safety alignment can be removed after release by fine-tuning~\citep{qi2024finetuning} or by deleting a refusal direction~\citep{arditi2024refusal}; it is concentrated in the first output tokens~\citep{qi2025safety} and holds only within a local basin of the weights~\citep{peng2024navigating}. Existing defenses each target a specific perturbation model: representation drift~\citep{huang2024vaccine,liu2024targeted}, harmful-representation suppression~\citep{rosati2024repnoise}, simulated or meta-learned attacks~\citep{huang2025booster,tamirisa2025tar}, constrained user fine-tuning~\citep{huang2024lisa,li2025salora}, and non-fine-tunable learning~\citep{deng2024sophon}. Such protection often fails when the attacker changes the learning rate, optimizer, data or objective~\citep{qi2025evaluating,zloczower2026one}.

\paragraph{Self-destructive alignment.}
Self-destructive methods instead make safety removal degrade benign capability~\citep{henderson2023self}, by degrading harmful answers~\citep{chen2025sdd}, embedding a collapse trap~\citep{yi2026ctrap}, or targeting harmful RL fine-tuning~\citep{feng2026tokenbuncher}. Closest to our work, SEAM~\citep{wang2026seam} anti-aligns the mean harmful and benign gradients and amplifies the effect with gradient ascent on harmful data. SEAM shapes the relation between two gradient directions of harmful adaptation; CBG instead penalizes the most collateral-efficient safety-changing directions over an admissible update space, regardless of their sign or loss. In a limiting case CBG reduces to a multi-direction analogue of SEAM's coupling (\cref{rem:seam}), and the two can be combined.

\paragraph{Geometric tools.}
CBG builds on established tools: Fisher geometry~\citep{amari1998natural,martens2020new,kunstner2019limitations}, generalized eigenproblems and sums of top eigenvalues~\citep{parlett1998symmetric,golub2013matrix,overton1993optimality}, and scale-invariant sharpness~\citep{dinh2017sharp,foret2021sharpness,kwon2021asam}; HarmAlign also applies spectral analysis to safety~\citep{rosati2026harmalign}. Our contribution is not these tools but the formulation of self-destructive alignment as a worst-direction collateral problem and an objective that keeps it trainable on minibatches.



\section{Preliminaries}
\label{sec:preliminaries}

\paragraph{Threat model.}
A defender releases a safety-aligned model $f_\theta$, obtained by further training an aligned model $f_{\theta_0}$. A white-box adversary with harmful data $\mathcal D_{\mathrm{atk}}=\{(x_l,y_l^h)\}_{l=1}^{N}$ attempts to remove safety behavior. Starting from $\theta$, a standard attack minimizes
\begin{equation}
\Lhfa(\theta')=\mathbb E_{(x,y^h)\sim\mathcal D_{\mathrm{atk}}}\,\ell\bigl(f_{\theta'}(x),y^h\bigr)
\label{eq:hfa}
\end{equation}
over $\theta'$, where $\ell$ is the token-level cross-entropy. Adaptive adversaries may change the objective (e.g., by adding a penalty that preserves benign behavior), the optimizer, the learning rate, the trainable parameters or the parameterization (full fine-tuning or LoRA), or may edit weights directly. The adversary succeeds if the modified model is both harmful and still useful.

\paragraph{Admissible local updates.}
We write an attack-induced change as $u=\theta'-\theta$ and do not assume that it follows the harmful gradient. Let $\U\subseteq\R^p$ be a linear subspace of admissible directions; by default $\U$ spans the weight matrices of all linear layers in the transformer blocks and excludes embeddings, normalization gains and the output head. The set $\mathcal V$ of updates a particular attacker can realize need not be a subspace: full fine-tuning of these matrices gives $\mathcal V=\U$, whereas rank-$r$ LoRA gives the matrices of rank at most $r$. Whenever $\mathcal V\subseteq\U$, the barrier defined in \cref{sec:formulation} is at least as large over $\mathcal V$ as over $\U$ (\cref{prop:monotone}), so it suffices to analyze $\U$. Updates with components outside $\U$ are not covered by this argument and are evaluated empirically. A direction in $\U$ that substantially changes safety-sensitive behavior at little benign functional cost is a \emph{low-collateral escape}.



\section{Methodology}
\label{sec:method}

CBG asks how cheaply, in benign functional terms, an admissible parameter change can move safety-sensitive behavior. We define this cost (\cref{sec:formulation}), make it well posed on minibatches (\cref{sec:practical}), and optimize it at LLM scale (\cref{sec:objective}).

% ---------------------------------------------------------
\subsection{Escape Geometry}
\label{sec:formulation}

\paragraph{Measurements.}
Safety data $\mathcal D_h=\{(x_i,y_i^s,y_i^h)\}_{i=1}^{m}$ pair each prompt with a safe and a harmful response; benign data $\mathcal D_b=\{(x_j,y_j)\}_{j=1}^{n}$ contain ordinary instructions. Each safety example contributes a token-normalized preference margin
\begin{equation}
\begin{aligned}
h_i(\theta)={}&\tfrac{1}{|y_i^h|}\log p_\theta(y_i^h\mid x_i)\\
&-\tfrac{1}{|y_i^s|}\log p_\theta(y_i^s\mid x_i),
\end{aligned}
\label{eq:safety_margin}
\end{equation}
which increases as the model moves toward the harmful response. Each benign example contributes the log-likelihood of tokens sampled from the model itself: under teacher forcing on $y_j$ we draw $\hat y_{j,t}\sim p_\theta(\cdot\mid x_j,y_{j,<t})$, fix the samples within a step, and set
\begin{equation}
F_{b,j}(\vartheta)=\textstyle\sum_{t}\log p_\vartheta(\hat y_{j,t}\mid x_j,y_{j,<t}).
\label{eq:mc_fisher}
\end{equation}
With $F_h=(h_i)_{i\le m}$, the $m+n$ weighted measurements $\Ft=(F_h/\sqrt m;\,F_b/\sqrt n)$ are indexed by $e$, and their Jacobian restricted to $\U$ is $C=[\Ch;\Cb]$ with rows $c_e$. For $u\in\U$,
\begin{equation}
Q_h(u)=\|\Ch u\|_2^2,\qquad Q_b(u)=\|\Cb u\|_2^2
\label{eq:metrics}
\end{equation}
are the first-order \emph{safety-sensitive displacement} and \emph{benign collateral} of $u$. $Q_h$ is sign-agnostic: it measures how strongly $u$ moves safety preferences, not in which direction. Because targets are sampled, $\mathbb E_{\hat y}[Q_b(u)]$ is the token-summed Fisher form, i.e.\ twice the second-order benign KL divergence between the current and the modified model (\cref{prop:mcfisher}). $Q_b$ therefore charges any change of the benign predictive distribution, which observed targets would not~\citep{kunstner2019limitations}.

\paragraph{Weakest escape and barrier.}
The attacker's most favorable local option is the direction with the largest displacement per unit of collateral.
\begin{definition}
\label{def:escape}
The \emph{escape efficiency} and the \emph{collateral barrier} at displacement $\gamma>0$ are
\begin{equation}
\begin{gathered}
\Resc=\sup_{u\in\U:\,Q_h(u)>0}\frac{Q_h(u)}{Q_b(u)},\\
\Barrier_\gamma=\inf_{u\in\U}\bigl\{\tfrac12Q_b(u):Q_h(u)\ge\gamma^2\bigr\},
\end{gathered}
\label{eq:rho_esc}
\end{equation}
with $c/0=+\infty$ for $c>0$ and $\sup\emptyset=0$. A maximizer is a \emph{weakest escape}.
\end{definition}
With the conventions $1/0=+\infty$ and $1/{+\infty}=0$,
\begin{equation}
\Barrier_\gamma=\frac{\gamma^2}{2\Resc}
\label{eq:barrier_identity}
\end{equation}
in every case, and $\Resc$ is finite, equal to the largest generalized eigenvalue of $(\Ch^\top\Ch,\Cb^\top\Cb)$ on the benign support of $\U$, exactly when $\U\cap\ker\Cb\subseteq\U\cap\ker\Ch$ (\cref{thm:barrier}). This is a standard generalized eigenproblem; it motivates the objective and is not a main result.

\paragraph{Finite-sample degeneracy.}
On a minibatch, $\operatorname{rank}\Cb\le n\ll\dim\U$, and the kernel condition holds only if every safety-margin gradient lies in the span of the $n$ benign gradients (\cref{prop:degenerate}). Generically it fails, so $\Resc=+\infty$ and $\Barrier_\gamma=0$: the benign sample leaves most directions unobserved, and the undamped objective reports a free escape through them.



\subsection{Trainable Geometry}
\label{sec:practical}

\paragraph{Relative coordinates.}
Unobserved directions need a movement cost, but transformers admit exact function-preserving diagonal rescalings (e.g., an MLP up-projection row against the matching down-projection column), and a Euclidean cost can be lowered by them without any change in behavior (\cref{prop:reparam}). Following adaptive sharpness~\citep{kwon2021asam}, we measure updates in relative units,
\begin{equation}
u=D_\theta\delta,\qquad D_\theta=\diag(\theta\odot\theta+\xi^2)^{1/2},
\label{eq:relative}
\end{equation}
so that $\|\delta\|_2^2=u^\top\Lam u$ with $\Lam=D_\theta^{-2}$. The relative Jacobian is $\widetilde C=CD_\theta=[\Cth;\Ctb]$, and its Gram $K=\widetilde C\widetilde C^\top=C\Lam^{-1}C^\top$, with blocks $K_{hh},K_{hb},K_{bb}$, is $(m+n)\times(m+n)$. For $\xi=0$ every quantity below is invariant under diagonal rescalings that preserve $F_h$ and $F_b$; the floor $\xi>0$ makes this approximate for $|\theta_i|\lesssim\xi$, and symmetries that mix coordinates are not covered.

\paragraph{Damping as soft projection.}
We charge movement $\mu\|\delta\|^2$ in addition to collateral. The damped efficiency is
\begin{equation}
\begin{aligned}
\rhoreg&=\max_{\delta\ne0}\frac{\|\Cth\delta\|^2}{\|\Ctb\delta\|^2+\mu\|\delta\|^2}
=\frac{\lambda_{\max}\bigl(\Cth\Pmu\Cth^{\top}\bigr)}{\mu},\\
\Pmu&=\mu\bigl(\Ctb^{\top}\Ctb+\mu I\bigr)^{-1},
\end{aligned}
\label{eq:damped}
\end{equation}
and the damped barrier $\Barrier^{(\mu)}_\gamma$, the least total cost $\tfrac12\bigl(Q_b(u)+\mu\,u^\top\Lam u\bigr)$ over $u\in\U$ with $Q_h(u)\ge\gamma^2$, equals $\gamma^2/(2\rhoreg)$ without any kernel condition; moreover $\rhoreg\uparrow\Resc$ as $\mu\downarrow0$ (\cref{prop:damped}). The matrix $0\preceq\Pmu\preceq I$ is a \emph{soft projector}: a direction with benign sensitivity $s^2$ is shrunk by $\mu/(s^2+\mu)$, and an unobserved direction passes unchanged. By the push-through identity, the projected safety Gram is a Schur complement of $K$,
\begin{equation}
\begin{aligned}
\Cth\Pmu\Cth^{\top}=\Smu:={}&K_{hh}\\
&-K_{hb}(\mu I_n+K_{bb})^{-1}K_{hb}^\top ,
\end{aligned}
\label{eq:reduced}
\end{equation}
i.e.\ the safety Gram with the part explained by benign sensitivity removed. We set $\mu=\kappa_\mu\tr(K_{bb})/n$, a multiple of the mean eigenvalue of the benign Gram, so $\kappa_\mu$ fixes the sensitivity below which a benign direction counts as unobserved.

\paragraph{Normalized top-$k$ efficiency.}
Minimizing $\rhoreg$ leaves two loopholes. \emph{Safety flattening}: $\rhoreg$ scales with $\|\Cth\|^2$, so shrinking safety sensitivity lowers it without making any escape costlier. \emph{Mode substitution}: suppressing only the leading escape lets the next one dominate. We therefore train the fraction of the top-$k$ relative safety sensitivity that survives the soft projection,
\begin{equation}
\begin{gathered}
\etak(\theta)=\frac{\sigma_{1:k}}{\nu_{1:k}}\in(0,1],\\
\sigma_j=\lambda_j(\Smu),\qquad \nu_j=\lambda_j(K_{hh}),
\end{gathered}
\label{eq:eta}
\end{equation}
where $\sigma_{1:k}=\sum_{j\le k}\sigma_j$, likewise $\nu_{1:k}$, and $k\le\operatorname{rank}\Cth\le m$. It is at most one because $\Smu\preceq K_{hh}$, and it needs only two $m\times m$ eigenproblems. It is unchanged by uniform safety flattening ($\Cth\mapsto t\Cth$) and by uniform benign inflation ($\Ctb\mapsto t\Ctb$, since $\mu$ scales with $K_{bb}$). For $k=1$, $1/\eta_1$ is the factor by which benign collateral raises the cost of the cheapest local escape over a movement-only cost (\cref{prop:eta}); as $\mu\downarrow0$, $\eta_1$ tends to a squared sine between the safety and benign gradient spans, a multi-direction, sign-agnostic analogue of SEAM's gradient coupling (\cref{rem:seam}).

The escapes follow from the same Gram. For $\Smu\phi_j=\sigma_j\phi_j$ with $\|\phi_j\|=1$, let
\begin{equation}
\begin{gathered}
z_j=(\mu I_n+K_{bb})^{-1}K_{hb}^\top\phi_j,\\
\psi_j=(\phi_j;\,-z_j)\in\R^{m+n}.
\end{gathered}
\label{eq:escape}
\end{equation}
Then $u_j=D_\theta\widetilde C^\top\psi_j/\sqrt{\mu\sigma_j}$ is the $j$-th weakest damped escape, with efficiency $\rho_j=\sigma_j/\mu$ and unit total cost (\cref{prop:reduced}).



\subsection{Objective and Optimization}
\label{sec:objective}

\paragraph{Objective.}
The defense minimizes
\begin{equation}
\widehat{\mathcal L}_{\mathrm{train}}=\Laln+\alpha\Lbgn+\beta\log\etak ,
\label{eq:training_objective}
\end{equation}
where $\Laln$ is the supervised loss of $y_i^s$ given $x_i$ and $\Lbgn$ is the token-level $\mathrm{KL}(p_{\theta_0}\Vert p_\theta)$ from the aligned model on benign data. The gradient flows through $\mu$: $\etak$ is invariant to benign inflation only because $\mu$ moves with $K_{bb}$, and detaching $\mu$ would reward inflating benign sensitivity (\cref{rem:mu}). The objective targets the collateral incurred per unit of safety change, not the amount of safety change an attack must achieve; it is therefore complementary to terms that increase the latter, such as SEAM's unlearning loss (\cref{prop:path}), and we evaluate the combination as CBG+UL. Minibatch spectral estimates are biased (\cref{prop:jensen}), so we use many short benign sequences ($n\ge4m$). The objective shapes local geometry; it is not a global robustness guarantee.

\paragraph{Gradient through the Gram.}
The objective depends on $\theta$ only through $K$, so its gradient is the pull-back of one small matrix. Let $K_{hh}\phi'_j=\nu_j\phi'_j$ with $\|\phi'_j\|=1$, $\psi'_j=(\phi'_j;\,0)$ and $E_b=0_m\oplus I_n$. At a spectral gap ($\sigma_k>\sigma_{k+1}$, $\nu_k>\nu_{k+1}$), $G=\partial\log\etak/\partial K$ is
\begin{equation}
\begin{aligned}
G={}&\frac{1}{\sigma_{1:k}}\textstyle\sum_{j\le k}\bigl(\psi_j\psi_j^\top+\tfrac{\kappa_\mu}{n}\|z_j\|^2E_b\bigr)\\
&-\frac{1}{\nu_{1:k}}\textstyle\sum_{j\le k}\psi'_j\psi_j'^\top
\end{aligned}
\label{eq:G}
\end{equation}
(\cref{prop:gradient}), where the $E_b$ term is the entire effect of differentiating $\mu$. Pulling $G$ back through $K=C\Lam^{-1}C^\top$ gives
\begin{equation}
\begin{gathered}
\nabla_\theta\log\etak=2\textstyle\sum_{e}H_{\Ft_e}w_e+d_\Lambda,\\
w_e=\Lam^{-1}C^\top G\mathbf 1_e,\quad d_\Lambda=2\,\theta\odot\operatorname{diag}(C^\top GC)
\end{gathered}
\label{eq:gram_grad}
\end{equation}
(\cref{prop:gram}), where $H_{\Ft_e}$ is the Hessian of the $e$-th measurement and $d_\Lambda$ comes from the dependence of $\Lam$ on $\theta$. Without a gap, the same expressions are Clarke subgradients of the two spectral sums.

\paragraph{Curvature term.}
The first term of \cref{eq:gram_grad} is a sum of $m+n$ Hessian--vector products with different directions. We obtain it from ordinary gradients by a centered difference along one shared direction. For any random pair $a,b\in\R^{m+n}$ with $\mathbb E[ab^\top]=G$, set $\Omega=a^\top\Ft$ and $w=\Lam^{-1}C^\top b$. Then
\begin{equation}
\begin{aligned}
\nabla_\theta\log\etak={}&\mathbb E_{a,b}\Bigl[\tfrac{1}{\varepsilon}\bigl(\nabla\Omega(\theta+\varepsilon w)\\
&-\nabla\Omega(\theta-\varepsilon w)\bigr)\Bigr]+d_\Lambda+\mathcal O(\varepsilon^2)
\end{aligned}
\label{eq:outer}
\end{equation}
(\cref{prop:randomized}). We draw $a=V|\Gamma|^{1/2}s$ and $b=V\operatorname{sgn}(\Gamma)|\Gamma|^{1/2}s$ from the eigendecomposition $G=V\Gamma V^\top$ and Rademacher signs $s$. The estimate is unbiased up to $\mathcal O(\varepsilon^2)$ but has sign-draw variance; draws can be averaged at two passes each. A deterministic alternative perturbs each example $e$ along its own $w_e$, which removes this variance at a higher cost per pass (\cref{app:impl}); \cref{sec:ablation} compares the two.

\paragraph{Algorithm.}
\Cref{alg:cbg} lists one step. The two perturbed passes reuse the minibatch, the samples $\hat y$ and the signs, and their average also yields the base-loss gradient up to $\mathcal O(\varepsilon^2)$. The step $\varepsilon$ sets the root-mean-square relative perturbation $\|\Lam^{1/2}\varepsilon w\|/\sqrt{\dim\U}$ to a radius $r$, read off the Gram through $\|\Lam^{1/2}w\|^2=b^\top Kb$. The perturbation is applied to fp32 master weights~\citep{micikevicius2018mixed}, and $r$ is chosen against an exact Hessian--vector product on a small model so that the compute precision resolves it (\cref{app:impl}).

\begin{algorithm}[t]
\caption{One CBG training step}
\label{alg:cbg}
\small
\begin{algorithmic}[1]
\Require weights $\theta$ (fp32); safety batch $(x_i,y_i^s,y_i^h)_{i\le m}$; benign batch $(x_j,y_j)_{j\le n}$; $\alpha,\beta,\kappa_\mu,k,\xi,r$
\State Forward at $\theta$; sample $\hat y$ from benign logits; form $\Ft$; store layer inputs $X$
\State Backward of $\mathbf 1^\top\Ft$; store output gradients $\Delta$; accumulate $K$ from $(X,\Delta)$
\State $\mu\gets\kappa_\mu\tr(K_{bb})/n$; top-$k$ eigenpairs of $\Smu$ and $K_{hh}$; $\etak$ and $G$ (\cref{eq:reduced,eq:eta,eq:G})
\State Draw $s$; form $a,b$; one chunked pass over $(X,\Delta)$ gives $w=\Lam^{-1}C^\top b$ and $d_\Lambda$
\State $\varepsilon\gets r\sqrt{\dim\U/(b^\top Kb)}$;\ $g\gets0$
\For{$\tau\in\{+1,-1\}$}
  \State Forward at $\theta+\tau\varepsilon w$ with the same $\hat y$
  \State $g\mathrel{+}=\nabla\bigl[\tau\tfrac{\beta}{\varepsilon}\,a^\top\Ft+\tfrac12(\Laln+\alpha\Lbgn)\bigr]$
\EndFor
\State $\theta\gets\textsc{OptimizerStep}(\theta,\ g+\beta d_\Lambda)$
\end{algorithmic}
\end{algorithm}



\providecommand{\tbd}{--}

\section{Experiments}
\label{sec:experiments}

We evaluate CBG in the order in which the experiments are run. Hyperparameters are selected on a validation split only (\cref{sec:exp_selection}). The main comparison against SEAM, SDD and CTRAP uses six fixed scenarios (\cref{sec:exp_main}). All remaining experiments compare CBG with the undefended base model only (\cref{sec:exp_lr,sec:exp_variance,sec:exp_decoding,sec:adaptive,sec:orthogonal,sec:barrier_analysis,sec:ablation,sec:overhead}).



\subsection{Experimental Setting}
\label{sec:setting}

\paragraph{Models.}
All defenses start from Llama-2-7B-Chat (\texttt{meta-llama/Llama-2-7b-chat-hf}, revision pinned at data preparation). The same checkpoint is the frozen reference $\theta_0$ of $\Lbgn$ and the \emph{Base} model of every table.

\paragraph{Data.}
We follow SEAM's data recipe~\citep{wang2026seam}. The BeaverTails~\citep{ji2023beavertails} rows with refusals are filtered and shuffled with seed 42, then split into prompt-disjoint pools:
\begin{itemize}
  \item \textbf{Defense set $\mathcal D_h$}: 4{,}000 safety pairs $(x, y^h, y^s)$ with harmful responses $y^h$ and refusals $y^s$.
  \item \textbf{Benign set $\mathcal D_b$}: 4{,}000 Alpaca-cleaned samples~\citep{taori2023alpaca}.
  \item \textbf{Harmful attack set}: 1{,}000 pairs disjoint from $\mathcal D_h$.
  \item \textbf{Harmfulness test set}: 100 BeaverTails \texttt{30k\_test} prompts, after SEAM's category filter.
  \item \textbf{Validation attack set} (CBG only, \cref{sec:exp_selection}): 1{,}000 rows, prompt-disjoint from all of the above.
  \item \textbf{Geometry set $\mathcal D_{\mathrm{geo}}$}: 8 fixed batches of $m=8$ pairs and $n=32$ Alpaca rows, disjoint from $\mathcal D_h$ and $\mathcal D_b$.
\end{itemize}
Benign fine-tuning data (\cref{sec:exp_main}) are 1{,}000 further Alpaca-cleaned rows, disjoint from $\mathcal D_b$ and $\mathcal D_{\mathrm{geo}}$. Every text uses SEAM's format \texttt{Question: \{prompt\}\textbackslash nAnswer: \{response\}} and is truncated to 256 tokens.

\paragraph{CBG.}
\begin{itemize}
  \item \textbf{Measured set and batch:} $\U$ is the weight of every linear layer inside the transformer blocks (embeddings, norms and the output head are excluded). Each CBG step uses $m=8$ safety pairs and $n=32$ benign sequences.
  \item \textbf{Geometry hyperparameters:} $k=4$, $\kappa_\mu=1$, $\xi=10^{-3}$, $\alpha=1$; $\beta$ is selected on the validation split.
  \item \textbf{Curvature term:} computed by an exact Hessian--vector product in TF32. We check it against an FP32 Hessian--vector product and a central finite difference before training (\cref{app:impl}).
  \item \textbf{Lazy geometry:} the full CBG step runs every 4 optimizer steps, and its geometry gradient is reused on the intermediate steps.
  \item \textbf{Optimization:} full-parameter AdamW, learning rate $2\times10^{-5}$, cosine schedule with 10\% warmup, 500 steps, model split by layer over two H200 GPUs.
  \item \textbf{CBG+UL:} additionally minimizes SEAM's unlearning loss $\mathcal L_{\mathrm{ul}}$ (layer-wise, log-transformed) on $y^h$, with weight $\gamma_{\mathrm{ul}}$ (\cref{sec:objective}).
\end{itemize}

\paragraph{Baselines.}
In the main comparison (\cref{sec:exp_main}) we compare against three defenses:
\begin{itemize}
  \item SEAM~\citep{wang2026seam};
  \item SDD~\citep{chen2025sdd};
  \item CTRAP~\citep{yi2025ctrap}.
\end{itemize}
Each defense is trained from the same Llama-2-7B-Chat checkpoint with its released code and recommended hyperparameters. We report the defense data and training budget of each method in \cref{app:hparams}. All other experiments compare CBG with the undefended Base model only.

\paragraph{Attack protocol.}
Every attack starts from the clean (un-attacked) checkpoint of the evaluated model; attacks are never chained. Unless stated otherwise:
\begin{itemize}
  \item full-parameter supervised fine-tuning for one epoch, batch size 4 (250 steps on 1{,}000 samples);
  \item cosine schedule, 10\% warmup, weight decay $0.01$, no gradient clipping, BF16, seed 42.
\end{itemize}
To fit a single 48\,GB GPU we use paged 8-bit AdamW with gradient checkpointing, in place of SEAM's FP32-state AdamW. This setting is identical for every model in a table, so all comparisons are within one attack configuration.

\paragraph{Metrics.}
\begin{itemize}
  \item \textbf{Harmfulness score (HS):} the mean probability of the \texttt{harmful} label assigned by the BeaverTails DeBERTa classifier\footnote{\texttt{domenicrosati/deberta-v3-xsmall-beavertails-harmful-qa-classifier}.} to prompt plus response. Responses are generated greedily with at most 50 new tokens, and an empty response scores 0. We also report the fraction of responses with score $>0.5$.
  \item \textbf{Zero-shot score (ZS):} the mean zero-shot accuracy (\texttt{acc}) on TruthfulQA (MC1), MMLU, HellaSwag and ARC-Easy, using LM Evaluation Harness~\citep{gao2024harness}.
  \item \textbf{Compromised:} as in SEAM, an attacked model is \emph{compromised} if HS $>10\%$ while ZS $>30\%$. Low HS together with low ZS is \emph{self-destruction}. Low HS with intact ZS and a high attack loss is a failed attack, not a defense, so we report the attack's training loss with every attacked row.
  \item \textbf{Geometry:} we report $\etak$ and $\rhohat$, always together with $\mu$, $m$ and $n$, and the normalizers $\nu_{1:k}$ and $\tr K_{bb}$, all on $\mathcal D_{\mathrm{geo}}$. The minibatch $\Resc$ and $\Barrier_\gamma$ are not reported, since they are $+\infty$ and $0$ by construction (\cref{prop:degenerate}).
\end{itemize}



\subsection{Hyperparameter Selection}
\label{sec:exp_selection}

Before any test-set evaluation, each CBG configuration (value of $\beta$, and $\gamma_{\mathrm{ul}}$ for CBG+UL) passes two gates that use no test data:
\begin{enumerate}
  \item \textbf{Geometry gate} (during training, step 100): training stops unless the held-out $\log\etak$ on $\mathcal D_{\mathrm{geo}}$ has decreased by at least $0.5$ from its value at step 0.
  \item \textbf{Validation-attack gate} (after training): one harmful attack at learning rate $5\times10^{-5}$ on the validation attack set. Harmfulness is measured on the held-out validation prompts, and the configuration is rejected if HS $\geq 75\%$.
\end{enumerate}
Among the configurations that pass, we select the one with the lowest validation HS whose pre-attack ZS is within $2$ points of Base. \Cref{tab:selection} lists every configuration tried, including rejected ones.

\begin{table*}[t]
\centering
\small
\caption{CBG configurations on the validation split. $\Delta\log\etak$: change of held-out $\log\etak$ ($\mu$, $m=8$, $n=32$) at step 100 and at the end of training. Val.\ HS: harmfulness after the validation attack ($5\times10^{-5}$). Pre ZS: zero-shot score before any attack.}
\label{tab:selection}
\begin{tabular}{lcccccc}
\toprule
Variant & $\beta$ & $\gamma_{\mathrm{ul}}$ & $\Delta\log\etak$ (100 / end) & Val.\ HS$\downarrow$ & Pre ZS$\uparrow$ & Selected\\
\midrule
CBG    & \tbd & -- & \tbd\ / \tbd & \tbd & \tbd & \tbd\\
CBG    & \tbd & -- & \tbd\ / \tbd & \tbd & \tbd & \tbd\\
CBG+UL & \tbd & \tbd & \tbd\ / \tbd & \tbd & \tbd & \tbd\\
CBG+UL & \tbd & \tbd & \tbd\ / \tbd & \tbd & \tbd & \tbd\\
\bottomrule
\end{tabular}
\end{table*}



\subsection{Main Comparison with SEAM, SDD and CTRAP}
\label{sec:exp_main}

We compare all defenses in six scenarios (\cref{tab:scenarios}): the clean model, harmful fine-tuning at three learning rates, benign fine-tuning, and fine-tuning on poisoned data. Benign fine-tuning measures whether a defense degrades under ordinary use, where self-destruction is undesirable. The poisoned scenario mixes 100 harmful and 900 benign samples (10\% contamination), a realistic attack that hides harmful data in a benign task. Results are in \cref{tab:main}.

\begin{table*}[t]
\centering
\small
\caption{Scenarios of the main comparison. Harmful samples come from the attack set; benign samples are Alpaca-cleaned rows disjoint from all defense data. All fine-tuning follows the attack protocol of \cref{sec:setting}.}
\label{tab:scenarios}
\begin{tabular}{llcl}
\toprule
Case & Fine-tuning data & Learning rate & Purpose\\
\midrule
\texttt{clean}          & none                       & --                 & Safety and utility of the defended checkpoint\\
\texttt{harmful-lr2e-5} & 1{,}000 harmful pairs      & $2\times10^{-5}$  & Robustness to harmful fine-tuning\\
\texttt{harmful-lr5e-5} & 1{,}000 harmful pairs      & $5\times10^{-5}$  & Robustness at a larger learning rate\\
\texttt{harmful-lr1e-4} & 1{,}000 harmful pairs      & $1\times10^{-4}$  & Robustness at a high learning rate\\
\texttt{benign-lr2e-5}  & 1{,}000 Alpaca-cleaned     & $2\times10^{-5}$  & Effect of ordinary fine-tuning on safety and utility\\
\texttt{mixed10-lr5e-5} & 100 harmful + 900 benign   & $5\times10^{-5}$  & Robustness to 10\% poisoned fine-tuning data\\
\bottomrule
\end{tabular}
\end{table*}

\begin{table*}[t]
\centering
\caption{Main comparison on Llama-2-7B-Chat (seed 42). HS and ZS in \%. For every fine-tuned case, a defense succeeds if HS stays $\leq 10\%$ or ZS falls to $\leq 30\%$ (self-destruction), except under benign fine-tuning, where ZS should stay close to \texttt{clean}. Attack training losses are in \cref{app:attack_loss}.}
\label{tab:main}
\resizebox{\textwidth}{!}{%
\begin{tabular}{lcccccccccccc}
\toprule
& \multicolumn{2}{c}{\texttt{clean}}
& \multicolumn{2}{c}{\texttt{harmful-lr2e-5}}
& \multicolumn{2}{c}{\texttt{harmful-lr5e-5}}
& \multicolumn{2}{c}{\texttt{harmful-lr1e-4}}
& \multicolumn{2}{c}{\texttt{benign-lr2e-5}}
& \multicolumn{2}{c}{\texttt{mixed10-lr5e-5}}\\
\cmidrule(lr){2-3}\cmidrule(lr){4-5}\cmidrule(lr){6-7}\cmidrule(lr){8-9}\cmidrule(lr){10-11}\cmidrule(lr){12-13}
Model & HS$\downarrow$ & ZS$\uparrow$ & HS$\downarrow$ & ZS & HS$\downarrow$ & ZS & HS$\downarrow$ & ZS & HS$\downarrow$ & ZS$\uparrow$ & HS$\downarrow$ & ZS\\
\midrule
Base   & \tbd & \tbd & \tbd & \tbd & \tbd & \tbd & \tbd & \tbd & \tbd & \tbd & \tbd & \tbd\\
SEAM~\citep{wang2026seam} & \tbd & \tbd & \tbd & \tbd & \tbd & \tbd & \tbd & \tbd & \tbd & \tbd & \tbd & \tbd\\
SDD~\citep{chen2025sdd}   & \tbd & \tbd & \tbd & \tbd & \tbd & \tbd & \tbd & \tbd & \tbd & \tbd & \tbd & \tbd\\
CTRAP~\citep{yi2025ctrap} & \tbd & \tbd & \tbd & \tbd & \tbd & \tbd & \tbd & \tbd & \tbd & \tbd & \tbd & \tbd\\
\midrule
CBG    & \tbd & \tbd & \tbd & \tbd & \tbd & \tbd & \tbd & \tbd & \tbd & \tbd & \tbd & \tbd\\
CBG+UL & \tbd & \tbd & \tbd & \tbd & \tbd & \tbd & \tbd & \tbd & \tbd & \tbd & \tbd & \tbd\\
\bottomrule
\end{tabular}}
\end{table*}

\begin{table*}[t]
\centering
\small
\caption{Per-task zero-shot accuracy (\%) in the \texttt{clean} case, the components of ZS in \cref{tab:main}.}
\label{tab:clean_tasks}
\begin{tabular}{lccccc}
\toprule
Model & TruthfulQA & MMLU & HellaSwag & ARC-Easy & ZS\\
\midrule
Base   & \tbd & \tbd & \tbd & \tbd & \tbd\\
SEAM   & \tbd & \tbd & \tbd & \tbd & \tbd\\
SDD    & \tbd & \tbd & \tbd & \tbd & \tbd\\
CTRAP  & \tbd & \tbd & \tbd & \tbd & \tbd\\
CBG    & \tbd & \tbd & \tbd & \tbd & \tbd\\
CBG+UL & \tbd & \tbd & \tbd & \tbd & \tbd\\
\bottomrule
\end{tabular}
\end{table*}



\subsection{Full Learning-Rate Sweep}
\label{sec:exp_lr}

To match SEAM's main robustness experiment, we extend the harmful attack to its full learning-rate grid, adding $8\times10^{-5}$ and $2\times10^{-4}$ to the three rates of \cref{tab:main}. \Cref{tab:lr_sweep} compares CBG with Base and reports each attack's training loss. A defense that only slows fine-tuning would show a higher loss but the same final HS; one that self-destructs would show a sharp drop in ZS.

\begin{table*}[t]
\centering
\caption{Harmful attack learning-rate sweep (1{,}000 samples, 250 steps, seed 42). HS and ZS in \%; Loss is the mean attack training loss.}
\label{tab:lr_sweep}
\resizebox{\textwidth}{!}{%
\begin{tabular}{lccccccccccccccc}
\toprule
& \multicolumn{3}{c}{$2\times10^{-5}$}
& \multicolumn{3}{c}{$5\times10^{-5}$}
& \multicolumn{3}{c}{$8\times10^{-5}$}
& \multicolumn{3}{c}{$1\times10^{-4}$}
& \multicolumn{3}{c}{$2\times10^{-4}$}\\
\cmidrule(lr){2-4}\cmidrule(lr){5-7}\cmidrule(lr){8-10}\cmidrule(lr){11-13}\cmidrule(lr){14-16}
Model & HS$\downarrow$ & ZS & Loss & HS$\downarrow$ & ZS & Loss & HS$\downarrow$ & ZS & Loss & HS$\downarrow$ & ZS & Loss & HS$\downarrow$ & ZS & Loss\\
\midrule
Base   & \tbd & \tbd & \tbd & \tbd & \tbd & \tbd & \tbd & \tbd & \tbd & \tbd & \tbd & \tbd & \tbd & \tbd & \tbd\\
CBG    & \tbd & \tbd & \tbd & \tbd & \tbd & \tbd & \tbd & \tbd & \tbd & \tbd & \tbd & \tbd & \tbd & \tbd & \tbd\\
CBG+UL & \tbd & \tbd & \tbd & \tbd & \tbd & \tbd & \tbd & \tbd & \tbd & \tbd & \tbd & \tbd & \tbd & \tbd & \tbd\\
\bottomrule
\end{tabular}}
\end{table*}



\subsection{Variance across Attack Seeds}
\label{sec:exp_variance}

The results above use a single attack seed. We repeat the attacks at $2\times10^{-5}$ and $1\times10^{-4}$ with attack seeds 43 and 44, together with the seed-42 run, on fixed defended checkpoints. This measures variability of the attack only, not of defense training. \Cref{tab:variance} reports mean and standard deviation over the three seeds.

\begin{table}[t]
\centering
\small
\caption{HS and ZS (\%, mean $\pm$ std over attack seeds 42--44).}
\label{tab:variance}
\begin{tabular}{lcccc}
\toprule
& \multicolumn{2}{c}{$2\times10^{-5}$} & \multicolumn{2}{c}{$1\times10^{-4}$}\\
\cmidrule(lr){2-3}\cmidrule(lr){4-5}
Model & HS$\downarrow$ & ZS & HS$\downarrow$ & ZS\\
\midrule
Base & \tbd\,$\pm$\,\tbd & \tbd\,$\pm$\,\tbd & \tbd\,$\pm$\,\tbd & \tbd\,$\pm$\,\tbd\\
CBG  & \tbd\,$\pm$\,\tbd & \tbd\,$\pm$\,\tbd & \tbd\,$\pm$\,\tbd & \tbd\,$\pm$\,\tbd\\
\bottomrule
\end{tabular}
\end{table}




\subsection{Sensitivity to Decoding}
\label{sec:exp_decoding}

HS is measured with greedy decoding throughout. To check that conclusions do not depend on this choice, we re-measure HS at $1\times10^{-4}$ with the checkpoint's default sampling configuration (Llama-2-Chat: temperature $0.6$, top-$p$ $0.9$), the setting of SEAM's evaluation script (\cref{tab:decoding}).

\begin{table}[t]
\centering
\small
\caption{HS (\%) under greedy and sampled decoding, before and after the $1\times10^{-4}$ attack.}
\label{tab:decoding}
\begin{tabular}{lcccc}
\toprule
& \multicolumn{2}{c}{Pre-attack} & \multicolumn{2}{c}{After $1\times10^{-4}$}\\
\cmidrule(lr){2-3}\cmidrule(lr){4-5}
Model & Greedy & Sampled & Greedy & Sampled\\
\midrule
Base & \tbd & \tbd & \tbd & \tbd\\
CBG  & \tbd & \tbd & \tbd & \tbd\\
\bottomrule
\end{tabular}
\end{table}




\subsection{Attack Variants}
\label{sec:exp_variants}

We next vary the attack itself, keeping learning rates of $5\times10^{-5}$ and $1\times10^{-4}$:
\begin{itemize}
  \item \textbf{Data size:} 10{,}000 harmful samples instead of 1{,}000.
  \item \textbf{Parameter-efficient fine-tuning:} LoRA~\citep{hu2022lora} with rank 8, $\alpha=16$, on the query/value projections and on all linear modules.
  \item \textbf{Optimizer:} SGD instead of AdamW.
  \item \textbf{Long training at small learning rates:} up to 50K steps at $10^{-6}$ and $5\times10^{-6}$.
  \item \textbf{Frozen layers:} fine-tuning with the intermediate layers frozen.
\end{itemize}
For each variant, \cref{tab:variants} reports HS, ZS and whether the attacked model is compromised.

\begin{table*}[t]
\centering
\small
\caption{Attack variants (HS / ZS in \%). Comp.: Yes if the attacked model is compromised (HS $>10\%$ and ZS $>30\%$), otherwise No.}
\label{tab:variants}
\begin{tabular}{lcccc}
\toprule
& \multicolumn{2}{c}{Base} & \multicolumn{2}{c}{CBG}\\
\cmidrule(lr){2-3}\cmidrule(lr){4-5}
Attack & HS / ZS & Comp. & HS / ZS & Comp.\\
\midrule
10K samples, $5\times10^{-5}$          & \tbd\ / \tbd & \tbd & \tbd\ / \tbd & \tbd\\
10K samples, $1\times10^{-4}$          & \tbd\ / \tbd & \tbd & \tbd\ / \tbd & \tbd\\
LoRA (q,v), $1\times10^{-4}$           & \tbd\ / \tbd & \tbd & \tbd\ / \tbd & \tbd\\
LoRA (all), $1\times10^{-4}$           & \tbd\ / \tbd & \tbd & \tbd\ / \tbd & \tbd\\
SGD, $1\times10^{-4}$                  & \tbd\ / \tbd & \tbd & \tbd\ / \tbd & \tbd\\
50K steps, $10^{-6}$                   & \tbd\ / \tbd & \tbd & \tbd\ / \tbd & \tbd\\
Frozen middle layers, $1\times10^{-4}$ & \tbd\ / \tbd & \tbd & \tbd\ / \tbd & \tbd\\
\bottomrule
\end{tabular}
\end{table*}




\subsection{Escape-Seeking Attacks}
\label{sec:adaptive}

An adaptive attacker who knows CBG can target low-collateral directions instead of following the plain harmful gradient. We evaluate four such attacks against CBG and Base (\cref{tab:adaptive}):
\begin{enumerate}
  \item \textbf{Utility-preserving:} adds $\lambda\,\mathrm{KL}(p_{\theta_0}\Vert p_\theta)$ on benign prompts to $\Lhfa$, with $\lambda\in\{0.01,0.05,0.1,0.2\}$ and $\theta_0$ the pre-attack defended model.
  \item \textbf{Projected:} removes from each attack gradient its component in the span of per-example benign gradients. This is the direction family that CBG penalizes as $\mu\to0$ (\cref{prop:degenerate}).
  \item \textbf{Reverse-CBG:} adds $-\beta'\log\etak$ to the attack loss, mirroring SEAM's trap reversal.
  \item \textbf{Renormalization:} a function-preserving rescaling of each SwiGLU block before fine-tuning (\cref{prop:reparam}).
\end{enumerate}
We also add Gaussian noise to the harmful gradient ($\sigma\in\{0.1,0.5,1\}$ relative to its RMS).

\begin{table*}[t]
\centering
\small
\caption{Escape-seeking attacks at $5\times10^{-5}$ (HS / ZS in \%).}
\label{tab:adaptive}
\begin{tabular}{lcc}
\toprule
Attack & Base & CBG\\
\midrule
Utility-preserving ($\lambda=0.1$) & \tbd\ / \tbd & \tbd\ / \tbd\\
Projected                          & \tbd\ / \tbd & \tbd\ / \tbd\\
Reverse-CBG                        & -- & \tbd\ / \tbd\\
Renormalization                    & \tbd\ / \tbd & \tbd\ / \tbd\\
Gradient noise ($\sigma=0.5$)      & \tbd\ / \tbd & \tbd\ / \tbd\\
\bottomrule
\end{tabular}
\end{table*}




\subsection{Orthogonalization}
\label{sec:orthogonal}

We apply refusal-direction orthogonalization~\citep{arditi2024refusal} to Base and CBG and report attack success rate (ASR) and ZS. This edit is not a fine-tuning step and lies outside our theory, so the experiment tests transfer only.

% ----------------------------------------------------------------------------
\subsection{Geometry Analysis}
\label{sec:barrier_analysis}

\paragraph{Training dynamics.}
We track the held-out $\log\etak$, $\nu_{1:k}$, $\tr K_{bb}$, $\Lbgn$ and $\Laln$ every 50 steps (\cref{fig:geometry}). Together they test whether $\etak$ falls through benign coupling, as intended, rather than through flattening, uniform inflation of benign gradients, or degradation of the model (rising $\Lbgn$).

\paragraph{Does geometry predict outcomes?}
Across CBG configurations (\cref{tab:selection}), checkpoints and attack learning rates, we correlate the pre-attack $\etak$ with the ZS drop at matched HS increase. We also measure the escape ratio at attack checkpoints, which tests the path-bound assumption of \cref{prop:path}.

\begin{table*}[t]
\centering
\fbox{\parbox{0.9\linewidth}{\centering\vspace{6em}Held-out $\log\etak$, $\nu_{1:k}$, $\tr K_{bb}$, $\Lbgn$ and $\Laln$ during CBG training.\vspace{6em}}}
\caption{Geometry and loss trajectories of CBG on $\mathcal D_{\mathrm{geo}}$ ($\mu$, $m=8$, $n=32$).}
\label{fig:geometry}
\end{table*}

% ----------------------------------------------------------------------------
\subsection{Ablations}
\label{sec:ablation}

Starting from the selected configuration, we change one component at a time and report held-out $\Delta\log\etak$, pre-attack ZS, and HS/ZS after the $5\times10^{-5}$ attack (\cref{tab:ablation}). The ablations are:
\begin{enumerate}
  \item $\log\etak$ versus raw $\rhohat$;
  \item $\Lam$ versus Euclidean damping;
  \item $k\in\{1,2,4,8\}$;
  \item sampled versus empirical Fisher;
  \item $n\in\{8,32,128\}$ at a fixed token budget;
  \item exact solve versus Rayleigh ascent with $K\in\{1,3,5\}$ iterations;
  \item exact Hessian--vector product versus finite differences, at several radii $r$ and in FP32 versus BF16;
  \item removing $\Laln$ or $\Lbgn$;
  \item detached versus differentiated $\mu$;
  \item one versus several sign draws, and the shared direction versus per-example perturbations;
  \item lazy period $1$ versus $4$;
  \item CBG versus CBG+UL versus UL alone ($\beta=0$).
\end{enumerate}
The last comparison isolates the contribution of the geometry term from that of the unlearning loss.

\begin{table*}[t]
\centering
\small
\caption{Ablations. Attack: $5\times10^{-5}$, HS / ZS in \%.}
\label{tab:ablation}
\begin{tabular}{lccc}
\toprule
Variant & $\Delta\log\etak$ & Pre ZS$\uparrow$ & Attack HS / ZS\\
\midrule
CBG (selected)              & \tbd & \tbd & \tbd\ / \tbd\\
UL only ($\beta=0$)         & \tbd & \tbd & \tbd\ / \tbd\\
CBG+UL                      & \tbd & \tbd & \tbd\ / \tbd\\
Raw $\rhohat$               & \tbd & \tbd & \tbd\ / \tbd\\
Euclidean damping           & \tbd & \tbd & \tbd\ / \tbd\\
Empirical Fisher            & \tbd & \tbd & \tbd\ / \tbd\\
Detached $\mu$              & \tbd & \tbd & \tbd\ / \tbd\\
Without $\Lbgn$             & \tbd & \tbd & \tbd\ / \tbd\\
Lazy period 1               & \tbd & \tbd & \tbd\ / \tbd\\
\bottomrule
\end{tabular}
\end{table*}

% ----------------------------------------------------------------------------
\subsection{Overhead}
\label{sec:overhead}

We report wall-clock time per step, peak GPU memory and total training time of CBG on the same hardware, broken down into center pass, Gram contraction, direction pass and curvature pass (\cref{app:impl}). For reference we also list the training time of SEAM, SDD and CTRAP as run in \cref{sec:exp_main}.

\begin{table*}[t]
\centering
\small
\caption{Training cost on 2$\times$H200 (500 steps).}
\label{tab:overhead}
\begin{tabular}{lccc}
\toprule
Method & Time / step (s) & Peak memory (GiB) & Total (h)\\
\midrule
SEAM   & \tbd & \tbd & \tbd\\
SDD    & \tbd & \tbd & \tbd\\
CTRAP  & \tbd & \tbd & \tbd\\
CBG    & \tbd & \tbd & \tbd\\
CBG+UL & \tbd & \tbd & \tbd\\
\bottomrule
\end{tabular}
\end{table*}





\section{Conclusion}
\label{sec:conclusion}

Self-destructive defenses make a chosen harmful adaptation costly, but an adaptive attacker can look for another route that removes safety at little benign cost. We proposed the no-cheap-escape principle, which makes the cheapest such route the defense target, and Collateral Barrier Geometry, which identifies the most collateral-efficient local directions out of safety and trains the model so that they incur benign damage. The broader takeaway is that a self-destructive defense should be judged by its weakest escape rather than by the attack it was trained against.

\section*{Limitations}

CBG's guarantees are local: they concern a first-order quadratic model at the released parameters and imply nothing about long fine-tuning trajectories, non-local edits such as weight orthogonalization, or updates outside the admissible space $\U$. The geometry is defined by the sampled safety and benign data, so behavior those samples do not cover is unmeasured. It is also estimated from small minibatches, which makes the trained quantity $\etak$ a biased, damping-dependent estimate with a noisy gradient. The normalization in $\etak$ rules out uniform safety flattening and uniform benign inflation only, and invariance holds only for diagonal rescaling symmetries. Each CBG step needs several forward and backward passes and stores per-layer inputs and output gradients (\cref{app:impl}). Our evaluation covers models of at most 8B parameters and English safety data.

\bibliographystyle{acl_natbib}
\bibliography{sn-bibliography}

% =========================================================
% Appendix
% =========================================================
\appendix

\section{Theoretical Analysis of Collateral Barrier Geometry}
\label{app:cbg_theory}

Throughout, $\U\subseteq\R^p$ is a nonzero linear subspace, $\Gh,\Gb$ are symmetric positive semidefinite, and $Q_h(u)=u^\top\Gh u$, $Q_b(u)=u^\top\Gb u$. Every quadratic form satisfies $Q(-u)=Q(u)$; positive semidefiniteness gives $Q\ge0$ and $u^\top Gu=0\Leftrightarrow Gu=0$. We write
\begin{equation*}
\begin{gathered}
\Nb=\U\cap\ker\Gb,\qquad \Nh=\U\cap\ker\Gh,\\
\mathcal S=\U\cap\Nb^\perp .
\end{gathered}
\end{equation*}
Standard facts are cited; results without a citation are ours.

% ---------------------------------------------------------
\subsection{Exact Semidefinite Collateral Geometry}
\label{app:cbg_exact_geometry}

\begin{proposition}[Zero-collateral escape]
\label{prop:zero}
If $\Nb\not\subseteq\Nh$, then for every $\gamma>0$ there is $u_\gamma\in\U$ with $Q_h(u_\gamma)=\gamma^2$ and $Q_b(u_\gamma)=0$.
\end{proposition}
\begin{proof}
Take $v\in\Nb\setminus\Nh$ and $u_\gamma=\gamma v/\sqrt{Q_h(v)}$.
\end{proof}

Under $\Nb\subseteq\Nh$, every $u\in\U$ decomposes uniquely as $u=s+n$ with $s\in\mathcal S$ and $n\in\Nb$, and then $Q_b(u)=Q_b(s)$ and $Q_h(u)=Q_h(s)$; moreover $\Gb$ is positive definite on $\mathcal S$. With $P_{\mathcal S}$ an orthonormal basis of $\mathcal S$, let
\begin{equation*}
A=P_{\mathcal S}^\top\Gh P_{\mathcal S}\succeq0,\qquad B=P_{\mathcal S}^\top\Gb P_{\mathcal S}\succ0,
\end{equation*}
and define $\Resc=\lambda_{\max}(A,B)$ when $Q_h\not\equiv0$ on $\U$.

\begin{theorem}[Semidefinite local collateral barrier]
\label{thm:barrier}
Fix $\gamma>0$.
(i)~If $Q_h\equiv0$ on $\U$, then $\Resc=0$ and $\Barrier_\gamma=+\infty$.
(ii)~If $\Nb\not\subseteq\Nh$, then $\Resc=+\infty$ and $\Barrier_\gamma=0$.
(iii)~Otherwise $0<\Resc<\infty$ and $\Barrier_\gamma=\gamma^2/(2\Resc)$, attained by a scaled dominant generalized eigenvector.
\end{theorem}
\begin{proof}
(i)~The constraint is infeasible. (ii)~By \cref{prop:zero}. (iii)~Some $u\in\U$ has $Q_h(u)=Q_h(s)>0$, so $A\ne0$ and $\Resc>0$; finiteness follows from $B\succ0$. By the decomposition the problem reduces to $\mathcal S$. For feasible $v$,
\begin{equation*}
\gamma^2\le v^\top Av\le\Resc\,v^\top Bv,
\end{equation*}
so $\tfrac12v^\top Bv\ge\gamma^2/(2\Resc)$. If $Av^\star=\Resc Bv^\star$ and $t=\gamma/\sqrt{{v^\star}^\top Av^\star}$, then $tv^\star$ is feasible with cost $\gamma^2/(2\Resc)$.
\end{proof}

\begin{remark}[The $+\infty$ convention is the supremum]
Let $\Nb\not\subseteq\Nh$ and $Q_b\not\equiv0$ on $\U$, and pick $v\in\Nb\setminus\Nh$ and $w\in\U$ with $Q_b(w)>0$. Since $\Gb v=0$, $Q_b(v+\epsilon w)=\epsilon^2Q_b(w)$ while $Q_h(v+\epsilon w)\to Q_h(v)>0$, so the ratio over $\{Q_b>0\}$ is unbounded. The convention is needed only when $Q_b\equiv0$ on $\U$.
\end{remark}

The theorem is exact for the specified quadratic model. It does not assert that $Q_h(u)\ge\gamma^2$ is equivalent to a successful jailbreak.

% ---------------------------------------------------------
\subsection{Admissible Sets}
\label{app:admissible}

\begin{proposition}[Monotonicity in the admissible set]
\label{prop:monotone}
Let $\mathcal V\subseteq\U$ be a cone ($tu\in\mathcal V$ for all $t>0$, $u\in\mathcal V$), and define
\begin{equation*}
\Barrier_\gamma(\mathcal V)=\inf\bigl\{\tfrac12Q_b(u):u\in\mathcal V,\ Q_h(u)\ge\gamma^2\bigr\},
\end{equation*}
and $\Barrier^{(\mu)}_\gamma(\mathcal V)$ analogously. Then $\Barrier_\gamma(\mathcal V)\ge\Barrier_\gamma(\U)$ and $\Barrier^{(\mu)}_\gamma(\mathcal V)\ge\Barrier^{(\mu)}_\gamma(\U)$. In particular, the modeled barrier against rank-$r$ LoRA on matrices contained in $\U$ is at least the barrier over $\U$.
\end{proposition}
\begin{proof}
An infimum over a subset is at least the infimum over the set, and the rank-$\le r$ updates of a matrix form a cone.
\end{proof}

If $u\mapsto Tu$ is an injective linear parameterization with range $\U$, the undamped efficiency equals $\lambda_{\max}(T^\top\Gh T,T^\top\Gb T)$ whenever the second matrix is positive definite, and it does not depend on $T$. The damped efficiency does, unless the damping transforms with $T$ (\cref{prop:reparam}).

% ---------------------------------------------------------
\subsection{Damping, Reparameterization, and Finite Samples}
\label{app:damping}

Let $\Lambda\succ0$ on $\U$ and $\mu>0$, and define $\rhoreg$ and $\Barrier^{(\mu)}_\gamma$ as in \cref{eq:damped}.

\begin{proposition}[Damped barrier and sandwich]
\label{prop:damped}
(i)~If $Q_h\not\equiv0$ on $\U$, then $\Barrier^{(\mu)}_\gamma=\gamma^2/(2\rhoreg)$; otherwise $\rhoreg=0$ and $\Barrier^{(\mu)}_\gamma=+\infty$.
(ii)~$\rhoreg\le\Resc$ always, strictly if $0<\Resc<\infty$.
(iii)~If $\Nb\subseteq\Nh$ and $Q_h\not\equiv0$ on $\U$, let
\begin{equation*}
\kappa_\Lambda=\min_{s\in\mathcal S\setminus\{0\}}\frac{Q_b(s)}{s^\top\Lambda s}>0 .
\end{equation*}
Then
\begin{equation}
\begin{gathered}
\rhoreg\le\Resc\le\Bigl(1+\frac{\mu}{\kappa_\Lambda}\Bigr)\rhoreg,\\
\frac{\Barrier^{(\mu)}_\gamma}{1+\mu/\kappa_\Lambda}\le\Barrier_\gamma\le\Barrier^{(\mu)}_\gamma .
\end{gathered}
\label{eq:app_sandwich}
\end{equation}
(iv)~$\rhoreg$ is nonincreasing in $\mu$, and $\rhoreg\to\Resc$ as $\mu\downarrow0$, including when $\Resc=+\infty$.
\end{proposition}
\begin{proof}
(i)~$\Gb+\mu\Lambda$ is positive definite on $\U$; apply the argument of \cref{thm:barrier}(iii).
(ii)~If $Q_b(u)>0$ and $Q_h(u)>0$, then
\begin{equation*}
\frac{Q_h(u)}{Q_b(u)+\mu u^\top\Lambda u}<\frac{Q_h(u)}{Q_b(u)}\le\Resc .
\end{equation*}
If $Q_b(u)=0$ and $\Resc<\infty$, then $Q_h(u)=0$. For strictness, the damped maximizer has $Q_h>0$, hence $Q_b>0$, and its damped ratio lies strictly below its undamped ratio.
(iii)~For $u=s+n$ the undamped ratio equals $Q_h(s)/Q_b(s)$, and $s^\top\Lambda s\le Q_b(s)/\kappa_\Lambda$ gives
\begin{equation*}
\frac{Q_h(s)}{Q_b(s)+\mu s^\top\Lambda s}\ge\frac{Q_h(s)}{Q_b(s)\,(1+\mu/\kappa_\Lambda)} .
\end{equation*}
Taking suprema gives the upper bound on $\Resc$ in \cref{eq:app_sandwich}, and (i) gives the barrier form.
(iv)~Monotonicity holds pointwise in $u$. If $\Resc<\infty$, (iii) gives convergence. If $\Resc=\infty$, take $v\in\Nb\setminus\Nh$; then $\rhoreg\ge Q_h(v)/(\mu\,v^\top\Lambda v)\to\infty$.
\end{proof}

\begin{proposition}[Reparameterization]
\label{prop:reparam}
Let $\psi$ be a diffeomorphism of parameter space with $F_h\circ\psi=F_h$ and $F_b\circ\psi=F_b$, and let $D=D\psi(\theta)$.
(i)~$\Ch(\psi(\theta))D=\Ch(\theta)$ and $\Cb(\psi(\theta))D=\Cb(\theta)$, so $\Resc$ and $\Barrier_\gamma$ computed at $\psi(\theta)$ over $D\U$ equal those at $\theta$ over $\U$.
(ii)~The damped quantities are invariant if and only if $D^\top\Lambda_{\psi(\theta)}D=\Lam$ on $\U$. For a diagonal symmetry $\psi_c(\theta)=\diag(c)\theta$ with $c>0$ and a coordinate subspace $\U$, the metric $\Lam=\diag(\theta_i^{-2})$ satisfies this, whereas $\Lam=I$ does so only if $c=\mathbf 1$ on $\U$.
(iii)~If $f\circ\psi_c=f$ for a one-parameter group $\psi_c$, then
\begin{equation*}
\bigl\langle\nabla f(\theta),\,\partial_c\psi_c(\theta)|_{c=1}\bigr\rangle=0 ;
\end{equation*}
a damped objective that is not invariant generically has a nonzero derivative along the orbit, so gradient descent can lower it by a function-preserving move.
\end{proposition}
\begin{proof}
(i)~Differentiate $\Ft\circ\psi=\Ft$ by the chain rule and substitute $u'=Du$.
(ii)~The damping term becomes $\mu u^\top D^\top\Lambda_{\psi(\theta)}Du$. For $D=\diag(c)$ and $\Lambda_{\psi(\theta)}=\diag((c_i\theta_i)^{-2})$ this equals $\mu\sum_iu_i^2/\theta_i^2$.
(iii)~Differentiate $f(\psi_c(\theta))=f(\theta)$ at $c=1$.
\end{proof}
Diagonal symmetries of pre-norm transformers include: a normalization gain against the corresponding input columns of the following projection; an MLP up-projection row against the matching down-projection column (the gated activation is linear in the up-projection); matching query and key dimensions, with rotary pairs scaled equally; and value dimensions against the matching output-projection columns. General linear symmetries are not covered, and the floor $\xi>0$ makes invariance approximate for $|\theta_i|\lesssim\xi$. This mirrors the reparameterization critique of sharpness~\citep{dinh2017sharp} and its adaptive remedy~\citep{kwon2021asam}.

\begin{proposition}[Finite-sample degeneracy and small-damping limit]
\label{prop:degenerate}
Identify $\U$ with $\R^d$ by a coordinate selection, let $\Ch\in\R^{m\times d}$ and $\Cb\in\R^{n\times d}$ be the restricted weighted Jacobians, and write $\widetilde C=C\Lambda^{-1/2}$ for either of them.
(i)~$\Nb\subseteq\Nh$ if and only if $\row(\Ch)\subseteq\row(\Cb)$. Hence if $\operatorname{rank}[\Ch;\Cb]>\operatorname{rank}\Cb$, in particular if $m+n\le d$ and the $m+n$ rows are linearly independent, then $\Resc=+\infty$ and $\Barrier_\gamma=0$.
(ii)~Let $P_\perp$ be the orthogonal projector onto $\ker\widetilde C_b$. Then, for all $\mu>0$,
\begin{equation}
\begin{gathered}
\mu\,\rhoreg\ \ge\ \lambda_{\max}\bigl(\widetilde C_hP_\perp\widetilde C_h^\top\bigr),\\
\lim_{\mu\downarrow0}\mu\,\rhoreg=\lambda_{\max}\bigl(\widetilde C_hP_\perp\widetilde C_h^\top\bigr).
\end{gathered}
\label{eq:app_smallmu}
\end{equation}
\end{proposition}
\begin{proof}
(i)~$\ker(C^\top C)=\ker C$, and $\ker\Cb\subseteq\ker\Ch$ if and only if $\row\Ch\subseteq\row\Cb$.
(ii)~By \cref{prop:reduced},
\begin{equation*}
\rhoreg=\lambda_{\max}\bigl(\widetilde C_h(\widetilde C_b^\top\widetilde C_b+\mu I)^{-1}\widetilde C_h^\top\bigr).
\end{equation*}
If $\widetilde C_b^\top\widetilde C_b=\sum_i\sigma_i^2e_ie_i^\top$, then
\begin{equation*}
\begin{aligned}
\mu(\widetilde C_b^\top\widetilde C_b+\mu I)^{-1}&=\sum_{\sigma_i>0}\frac{\mu}{\sigma_i^2+\mu}\,e_ie_i^\top+P_\perp\\
&\succeq P_\perp ,
\end{aligned}
\end{equation*}
which is nondecreasing in $\mu$ and tends to $P_\perp$. The map $Y\mapsto\lambda_{\max}(XYX^\top)$ is monotone in the Loewner order and continuous.
\end{proof}

\begin{remark}[Relation to SEAM]
\label{rem:seam}
By \cref{eq:app_smallmu}, with small damping the finite-sample objective penalizes the largest energy of per-example safety-margin gradients outside the span of per-example benign gradients. SEAM's cosine loss~\citep{wang2026seam} drives one mean harmful gradient toward the negative of one mean benign gradient, i.e.\ into a one-dimensional benign span with a prescribed sign. CBG is a multi-direction, sign-agnostic analogue over every example in the batch: for $k=1$ and $\mu\downarrow0$,
\begin{equation*}
\eta_1\ \to\ \frac{\lambda_{\max}(\widetilde C_hP_\perp\widetilde C_h^\top)}{\lambda_{\max}(\widetilde C_h\widetilde C_h^\top)},
\end{equation*}
a generalized squared sine between the two spans. Moreover, the damped ratio of every $u\in\U$ is at most $\rhoreg$, including the restriction to $\U$ of the harmful gradient that SEAM targets.
\end{remark}

\begin{proposition}[Minibatch bias and consistency]
\label{prop:jensen}
Let $\Ghh$ and $\Gbh$ be independent, integrable random positive-semidefinite matrices with means $\bar{\mathcal G}_h$ and $\bar{\mathcal G}_b$, and let $f(A,B)=\lambda_{\max}(A,B+\mu\Lambda)$. Then
\begin{equation*}
\mathbb E\,f(\Ghh,\Gbh)\ \ge\ f(\bar{\mathcal G}_h,\bar{\mathcal G}_b).
\end{equation*}
If the minibatch metrics are averages of $m$ and $n$ i.i.d.\ terms, then $f(\Ghh,\Gbh)\to f(\bar{\mathcal G}_h,\bar{\mathcal G}_b)$ almost surely as $m,n\to\infty$.
\end{proposition}
\begin{proof}
For fixed $B$, $A\mapsto f$ is a pointwise maximum of linear functions, hence convex. For fixed $A\succeq0$, the map $B\mapsto(u^\top Au)/(u^\top(B+\mu\Lambda)u)$ is convex on the positive-definite cone, as a nonnegative multiple of the convex function $1/x$ composed with an affine map~\citep{boyd2004convex}, and so is its maximum over $u$. Jensen's inequality in $A$, conditionally on the independent $B$, and then in $B$ gives the bound. Consistency follows from the strong law of large numbers in finite dimension and continuity of $f$ for $\mu>0$; matrix Bernstein bounds give rates under bounded per-example gradients~\citep{tropp2012user}.
\end{proof}

\begin{proposition}[Sampled per-token Fisher]
\label{prop:mcfisher}
Fix $\theta$, a prompt $x$ and a reference $y$. Draw $\hat y_t\sim p_\theta(\cdot\mid x,y_{<t})$ independently over $t$, and let $s=\sum_t\nabla_\theta\log p_\theta(\hat y_t\mid x,y_{<t})$. Then $\mathbb E[ss^\top]=\sum_tF_t$, where $F_t$ is the Fisher information of $p_\theta(\cdot\mid x,y_{<t})$:
\begin{equation*}
F_t=\mathbb E_{v\sim p_\theta(\cdot\mid x,y_{<t})}\bigl[g_t(v)g_t(v)^\top\bigr],
\end{equation*}
with $g_t(v)=\nabla_\theta\log p_\theta(v\mid x,y_{<t})$. Moreover,
\begin{equation*}
\begin{aligned}
&\textstyle\sum_t D_{\mathrm{KL}}\bigl(p_\theta(\cdot\mid x,y_{<t})\,\Vert\,p_{\theta+u}(\cdot\mid x,y_{<t})\bigr)\\
&\quad=\tfrac12u^\top\Bigl(\textstyle\sum_tF_t\Bigr)u+o(\|u\|_2^2).
\end{aligned}
\end{equation*}
\end{proposition}
\begin{proof}
The score has zero mean under the model, so the cross terms $\mathbb E[s_ts_{t'}^\top]=\mathbb E[s_t]\mathbb E[s_{t'}]^\top$ vanish for $t\ne t'$ by independence. The KL expansion is the standard Fisher identity~\citep{amari1998natural,martens2020new}.
\end{proof}

% ---------------------------------------------------------
\subsection{Normalized Collateral Efficiency}
\label{app:eta}

\begin{proposition}[Properties of $\etak$]
\label{prop:eta}
Let $\Bmu=\Gbh+\mu\Lam$ and $1\le k\le\operatorname{rank}(\Ghh|_\U)$.
(i)~$\rho_j\le\nu_j/\mu$ for every $j$, hence $0<\etak\le1$.
(ii)~$\etak$ is unchanged under $\Ch\mapsto t\Ch$, under $(\Cb,\mu)\mapsto(t\Cb,t^2\mu)$, and under the diagonal symmetries of \cref{prop:reparam} with $\xi=0$.
(iii)~For $k=1$, let
\begin{equation*}
\Barrier^{(\mu,0)}_\gamma=\inf\bigl\{\tfrac12\mu\,u^\top\Lam u:u\in\U,\ Q_h(u)\ge\gamma^2\bigr\}
\end{equation*}
be the cost of the cheapest escape when benign collateral is ignored. Then $\Barrier^{(\mu)}_\gamma=\Barrier^{(\mu,0)}_\gamma/\eta_1$.
\end{proposition}
\begin{proof}
(i)~By the Courant--Fischer characterization for definite pencils~\citep{parlett1998symmetric},
\begin{equation*}
\rho_j=\max_{\dim\mathcal V=j}\ \min_{u\in\mathcal V\setminus\{0\}}\frac{u^\top\Ghh u}{u^\top\Bmu u},
\end{equation*}
and $\Bmu\succeq\mu\Lam$ together with $u^\top\Ghh u\ge0$ gives $\rho_j\le\nu_j/\mu$. Positivity holds because $\rho_1>0$ when $\Ghh|_\U\ne0$.
(ii)~In the first case numerator and denominator both scale by $t^2$. In the second, $\Bmu\mapsto t^2\Bmu$, so $\rho_j\mapsto\rho_j/t^2$ and $\mu\rho_j$ is unchanged. The third follows from \cref{prop:reparam}.
(iii)~Apply \cref{prop:damped}(i) with $\Gbh$ replaced by $0$: $\Barrier^{(\mu,0)}_\gamma=\gamma^2\mu/(2\nu_1)$, while $\Barrier^{(\mu)}_\gamma=\gamma^2/(2\rho_1)$.
\end{proof}

\begin{remark}[Detached versus differentiated $\mu$]
\label{rem:mu}
Property (ii) concerns values. The value of $\etak$ is invariant under $(\Cb,\mu)\mapsto(t\Cb,t^2\mu)$ only because $\mu=\kappa_\mu\tr(K_{bb})/n$ moves with the benign Gram. With $\mu$ detached, the gradient ignores this, and uniformly inflating benign sensitivity lowers $\etak$ (\cref{tab:numerical}). Since $\partial\Smu/\partial\mu=K_{hb}(\mu I_n+K_{bb})^{-2}K_{hb}^\top$,
\begin{equation*}
c_\mu:=\frac{\partial\log\etak}{\partial\mu}=\frac{\sum_{j\le k}\|z_j\|^2}{\sigma_{1:k}}\ \ge0,
\end{equation*}
and differentiating $\mu$ adds $(\kappa_\mu/n)\,c_\mu E_b$ to $\partial\log\etak/\partial K$, the $E_b$ term of \cref{eq:G}. Its coefficient is nonnegative, so it raises the objective when $\tr K_{bb}$ grows and penalizes benign inflation. In terms of the damped escapes, $c_\mu=1/\mu-\sum_j\rho_j\|D_\theta^{-1}u_j\|^2/\sum_j\rho_j$, which is nonnegative because $\mu\|D_\theta^{-1}u_j\|^2\le Q_b(u_j)+\mu\|D_\theta^{-1}u_j\|^2=1$.
\end{remark}

% ---------------------------------------------------------
\subsection{Exact Reduced Solve}
\label{app:reduced}

\begin{proposition}[Reduced system]
\label{prop:reduced}
Let $\widetilde C=CD_\theta=[\Cth;\Ctb]$, $K=\widetilde C\widetilde C^\top$ and $\Smu$ as in \cref{eq:reduced}.
(i)~$\Cth\Pmu\Cth^\top=\Smu$. Hence the nonzero generalized eigenvalues of $(\Ch^\top\Ch,\ \Cb^\top\Cb+\mu\Lam)$ on $\U$ are $\sigma_j/\mu$, where $\sigma_j$ runs over the positive eigenvalues of $\Smu$ with multiplicities, and those of $(\Ch^\top\Ch,\Lam)$ are the positive eigenvalues $\nu_j$ of $K_{hh}$.
(ii)~If $\Smu\phi=\sigma\phi$ with $\sigma>0$ and $\|\phi\|=1$, let $z=(\mu I_n+K_{bb})^{-1}K_{hb}^\top\phi$, $\psi=(\phi;-z)$ and $u=D_\theta\widetilde C^\top\psi/\sqrt{\mu\sigma}$. Then
\begin{equation*}
\begin{gathered}
\Ch^\top\Ch u=\tfrac{\sigma}{\mu}\bigl(\Cb^\top\Cb+\mu\Lam\bigr)u,\\
Q_b(u)+\mu\,u^\top\Lam u=1,\\
\Ch u=\sqrt{\sigma/\mu}\,\phi,\qquad \Cb u=\sqrt{\mu/\sigma}\,z .
\end{gathered}
\end{equation*}
(iii)~If $K_{hh}\phi'=\nu\phi'$ with $\nu>0$ and $\|\phi'\|=1$, then $v=\Lam^{-1}\Ch^\top\phi'/\sqrt\nu$ satisfies $\Ch^\top\Ch v=\nu\Lam v$ and $v^\top\Lam v=1$.
\end{proposition}
\begin{proof}
Work in relative coordinates $\delta=D_\theta^{-1}u$, where the pencil becomes $(\Cth^\top\Cth,B)$ with $B=\Ctb^\top\Ctb+\mu I=\mu\Pmu^{-1}$. (i)~The nonzero eigenvalues of $B^{-1}\Cth^\top\Cth$ equal those of $\Cth B^{-1}\Cth^\top=\mu^{-1}\Cth\Pmu\Cth^\top$~\citep{horn2012matrix}. The push-through identity
\begin{equation*}
(\mu I+\Ctb^\top\Ctb)^{-1}\Ctb^\top=\Ctb^\top(\mu I_n+K_{bb})^{-1}
\end{equation*}
gives $\Pmu=I-\Ctb^\top(\mu I_n+K_{bb})^{-1}\Ctb$ and hence $\Cth\Pmu\Cth^\top=\Smu$~\citep{golub2013matrix}. (ii)~By the same identity, $\delta:=D_\theta^{-1}u=\Pmu\Cth^\top\phi/\sqrt{\mu\sigma}$. Then $\Cth\delta=\Smu\phi/\sqrt{\mu\sigma}=\sqrt{\sigma/\mu}\,\phi$; $\Ctb\Pmu=\mu(\mu I_n+K_{bb})^{-1}\Ctb$ gives $\Ctb\delta=\sqrt{\mu/\sigma}\,z$; $B\delta=\mu\Cth^\top\phi/\sqrt{\mu\sigma}$, so $\Cth^\top\Cth\delta=\sqrt{\sigma/\mu}\,\Cth^\top\phi=(\sigma/\mu)B\delta$; and $\delta^\top B\delta=\phi^\top\Smu\phi/\sigma=1$. Returning to $u=D_\theta\delta$ gives the claims, since $\Ch u=\Cth\delta$, $\Cb u=\Ctb\delta$ and $u^\top\Lam u=\|\delta\|^2$. (iii)~is (ii) with $\Ctb=0$ and $\mu=1$.
\end{proof}

% ---------------------------------------------------------
\subsection{Gradient and Its Estimation}
\label{app:outer}

\paragraph{Assumptions.} $F_h,F_b$ are $C^4$ near $\theta$; samples, sampled targets and $W_h,W_b$ are fixed within a step; $\U$ is fixed; and $\xi>0$, so that $\Lam$ is smooth.

\begin{proposition}[Gram pull-back]
\label{prop:gram}
Let $f$ be differentiable at $K(\theta)=C(\theta)\Lam^{-1}C(\theta)^\top$, with symmetric gradient $G$, i.e.\ $df=\tr(G\,dK)$. Then
\begin{equation*}
\begin{aligned}
\nabla_\theta f\bigl(K(\theta)\bigr)={}&2\textstyle\sum_eH_{\Ft_e}w_e\\
&+2\,\theta\odot\operatorname{diag}(C^\top GC),
\end{aligned}
\end{equation*}
with $w_e=\Lam^{-1}C^\top G\mathbf 1_e$. Moreover $\operatorname{diag}(C^\top GC)=\Lam\sum_ec_e\odot w_e$ and $\|\Lam^{1/2}w_e\|^2=(GKG)_{ee}$.
\end{proposition}
\begin{proof}
Since $dc_e=H_{\Ft_e}d\theta$ and $d(\theta_i^2+\xi^2)=2\theta_i\,d\theta_i$,
\begin{equation*}
\begin{aligned}
dK_{ee'}={}&c_e^\top\Lam^{-1}H_{\Ft_{e'}}d\theta+(H_{\Ft_e}d\theta)^\top\Lam^{-1}c_{e'}\\
&+2\textstyle\sum_ic_{e,i}c_{e',i}\theta_i\,d\theta_i .
\end{aligned}
\end{equation*}
Summing against the symmetric $G$ and using the symmetry of each Hessian,
\begin{equation*}
\begin{aligned}
\tr(G\,dK)={}&2\textstyle\sum_{e}\bigl(H_{\Ft_e}\sum_{e'}G_{ee'}\Lam^{-1}c_{e'}\bigr)^\top d\theta\\
&+2\textstyle\sum_i\theta_i[C^\top GC]_{ii}\,d\theta_i ,
\end{aligned}
\end{equation*}
and $\sum_{e'}G_{ee'}\Lam^{-1}c_{e'}=w_e$. For the last two claims, $[C^\top GC]_{ii}=\sum_ec_{e,i}[\Lam w_e]_i$ and $\Lam^{1/2}w_e=\widetilde C^\top G\mathbf 1_e$.
\end{proof}

\begin{proposition}[Gradient of $\log\etak$]
\label{prop:gradient}
Let $\mu=\kappa_\mu\tr(K_{bb})/n$, and let $\phi_j,z_j,\psi_j$ and $\phi'_j,\psi'_j$ be as in \cref{eq:escape,eq:G}. If $\sigma_k>\sigma_{k+1}$ and $\nu_k>\nu_{k+1}$, then $\log\etak$ is differentiable in $K$ with gradient $G$ of \cref{eq:G}. With $\mu$ held fixed, the gradient is \cref{eq:G} without the $E_b$ term. Without the gap conditions, $\sigma_{1:k}$ and $\nu_{1:k}$ are locally Lipschitz in $K$, and the same expressions, built from any orthonormal basis of the top-$k$ eigenspaces, are Clarke subgradients.
\end{proposition}
\begin{proof}
By Ky Fan's principle, $\sigma_{1:k}=\max_{Y^\top Y=I_k}\tr(Y^\top\Smu Y)$; at a gap its differential is $\tr(\Phi_k^\top d\Smu\,\Phi_k)$ with $\Phi_k=[\phi_1,\dots,\phi_k]$~\citep{kato1995perturbation,magnus1985differentiating}. With $R=(\mu I_n+K_{bb})^{-1}$ and $Z=RK_{hb}^\top\Phi_k$,
\begin{equation*}
\begin{aligned}
d\Smu={}&dK_{hh}-dK_{hb}RK_{hb}^\top-K_{hb}R\,dK_{hb}^\top\\
&+K_{hb}R\,(dK_{bb}+d\mu\,I_n)\,RK_{hb}^\top,
\end{aligned}
\end{equation*}
so $\tr(\Phi_k^\top d\Smu\,\Phi_k)=\tr(\Psi_k\Psi_k^\top dK)+\|Z\|_F^2\,d\mu$ with $\Psi_k=[\Phi_k;-Z]$. Since $d\mu=(\kappa_\mu/n)\tr(E_b\,dK)$, this is $\tr\bigl([\Psi_k\Psi_k^\top+(\kappa_\mu/n)\|Z\|_F^2E_b]\,dK\bigr)$. Likewise $d\nu_{1:k}=\tr(\Psi'_k\Psi_k'^\top dK)$, and $d\log\etak=d\sigma_{1:k}/\sigma_{1:k}-d\nu_{1:k}/\nu_{1:k}$. Without gaps, each spectral sum is a maximum of smooth functions over a compact set, so it is locally Lipschitz and the gradients at maximizers generate its Clarke subdifferential~\citep[Sec.~2.8]{clarke1983optimization}; see also \citet{overton1993optimality}. Stochastic subgradient methods converge for such functions~\citep{davis2020stochastic}.
\end{proof}

\begin{proposition}[Curvature estimators]
\label{prop:randomized}
(i)~\emph{Shared direction.} If $a,b\in\R^{m+n}$ are random with $\mathbb E[ab^\top]=G$, $\Omega=a^\top\Ft$ and $w=\Lam^{-1}C^\top b$, then $\mathbb E[2H_\Omega w]=2\sum_eH_{\Ft_e}w_e$. The eigen-sign pair $a=V|\Gamma|^{1/2}s$, $b=V\operatorname{sgn}(\Gamma)|\Gamma|^{1/2}s$ with $G=V\Gamma V^\top$ and Rademacher $s$ satisfies $\mathbb E[ab^\top]=G$, and $\|\Lam^{1/2}w\|^2=b^\top Kb$.
(ii)~\emph{Per-example directions.} For radii $\varepsilon_e>0$,
\begin{equation*}
\begin{aligned}
&\textstyle\sum_e\frac{1}{\varepsilon_e}\bigl(\nabla\Ft_e(\theta+\varepsilon_ew_e)-\nabla\Ft_e(\theta-\varepsilon_ew_e)\bigr)\\
&\qquad=2\textstyle\sum_eH_{\Ft_e}w_e+\mathcal O\bigl(\max_e\varepsilon_e^2\bigr).
\end{aligned}
\end{equation*}
(iii)~For any fixed $\Omega$ and $w$,
\begin{equation}
\begin{aligned}
&\Bigl\|\tfrac{1}{\varepsilon}\bigl(\nabla\Omega(\theta+\varepsilon w)-\nabla\Omega(\theta-\varepsilon w)\bigr)-2H_\Omega w\Bigr\|\\
&\qquad\le\frac{\varepsilon^2}{3}\sup_{|\tau|\le\varepsilon}\bigl\|\nabla^4\Omega(\theta+\tau w)[w,w,w]\bigr\| .
\end{aligned}
\label{eq:app_fd}
\end{equation}
\end{proposition}
\begin{proof}
(i)~$H_\Omega w=\sum_{e,e'}a_eb_{e'}H_{\Ft_e}\Lam^{-1}c_{e'}$; taking expectations and using $\sum_{e'}G_{ee'}\Lam^{-1}c_{e'}=w_e$ gives the claim. For the eigen-sign pair, $\mathbb E[ab^\top]=V|\Gamma|^{1/2}\,\mathbb E[ss^\top]\operatorname{sgn}(\Gamma)|\Gamma|^{1/2}V^\top=G$, and $\Lam^{1/2}w=\widetilde C^\top b$. (ii) and (iii) are the centered-difference Taylor remainder in integral form~\citep{nocedal2006numerical}: even-order terms cancel, and the two third-order remainders each contribute at most $\tfrac{\varepsilon^3}{6}\sup\|\nabla^4\Omega[w,w,w]\|$ before division by $\varepsilon$, applied per example in (ii).
\end{proof}
The shared direction applies the Rademacher device of Hutchinson's trace estimator~\citep{hutchinson1989stochastic} to a sum of Hessian--vector products with different directions. It is unbiased up to $\mathcal O(\varepsilon^2)$ but has sign-draw variance (\cref{tab:numerical}); several draws can be averaged at two passes each. The per-example form has no sign variance, because examples do not interact in the forward pass and each can be evaluated at its own perturbation; with $\varepsilon_e=r\sqrt{\dim\U/(GKG)_{ee}}$, every example is perturbed to the same relative radius $r$. The Hessian--vector products could also be computed exactly by reverse-over-reverse differentiation~\citep{pearlmutter1994fast}, but that requires double-backward support in fused attention kernels, which the two-pass difference avoids.

\paragraph{Computing $d_\Lambda$.}
By \cref{prop:gram}, $\operatorname{diag}(C^\top GC)=\sum_ec_e\odot\bigl(\sum_{e'}G_{ee'}c_{e'}\bigr)$ in terms of the per-example gradients $c_e$. It is accumulated layer by layer in the chunked pass over the stored layer inputs and output gradients that also forms $w$ (\cref{app:impl}), so it requires no additional forward or backward pass.

\paragraph{Precision.}
In bf16 the relative spacing of a weight is $2^{-8}$~\citep{kalamkar2019study}, so coordinates of $\varepsilon w$ smaller than about $2^{-8}|\theta_i|$ are rounded away and the difference in \cref{eq:app_fd} becomes rounding noise (\cref{app:numerical}). Two roundings matter: storing $\theta\pm\varepsilon w$, and rounding inside each matrix multiplication (bf16 autocast rounds weights and activations to $2^{-8}$ relative, TF32 to about $2^{-11}$). Applying the shared perturbation to fp32 master weights removes the first. The per-example form applies $x_e(\theta\pm\varepsilon_ew_e)^\top=x_e\theta^\top\pm\varepsilon_e\,x_ew_e^\top$ as an additive branch, which also keeps the perturbation out of the weights. Neither removes the second, so $r$ must lie above the compute resolution and is set by comparing against an exact Hessian--vector product on a small model. The classical optimum for centered differences scales as the cube root of the working precision~\citep{nocedal2006numerical}.

\begin{proposition}[Finite-radius objective]
\label{prop:finite}
For fixed $u$ define
\begin{equation*}
a_\varepsilon=\frac{\|\Ft_h(\theta+\varepsilon u)-\Ft_h(\theta-\varepsilon u)\|^2}{4\varepsilon^2},
\end{equation*}
$b_\varepsilon$ analogously with $\Ft_b$ plus $\mu u^\top\Lam u$, and $R_\varepsilon=a_\varepsilon/b_\varepsilon$. Then $\nabla_\theta R_\varepsilon(\theta;u^\star)=\nabla_\theta\rhoreg+\mathcal O(\varepsilon^2)$ at an exact simple maximizer $u^\star$.
\end{proposition}
\begin{proof}
Write $\Delta_\varepsilon F=\Ft_h(\theta+\varepsilon u)-\Ft_h(\theta-\varepsilon u)$ and $\Delta_\varepsilon C=\Ch(\theta+\varepsilon u)-\Ch(\theta-\varepsilon u)$. Then $\Delta_\varepsilon F=2\varepsilon\Ch u+\mathcal O(\varepsilon^3)$ and $\Delta_\varepsilon C=2\varepsilon\,\partial\Ch[u]+\mathcal O(\varepsilon^3)$, so
\begin{equation*}
\begin{aligned}
\nabla a_\varepsilon&=\tfrac{1}{2\varepsilon^2}(\Delta_\varepsilon C)^\top\Delta_\varepsilon F\\
&=2\,\partial\Ch[u]^\top\Ch u+\mathcal O(\varepsilon^2)=\nabla a+\mathcal O(\varepsilon^2),
\end{aligned}
\end{equation*}
and likewise for $b_\varepsilon$. The envelope identity at $u^\star$ completes the proof.
\end{proof}
The gradient of $R_\varepsilon$ is exact for the finite-radius objective, which includes curvature within radius $\varepsilon$; it is an optional variant costing two extra forward passes.

% ---------------------------------------------------------
\subsection{From Local Geometry to Paths}
\label{app:path}

\begin{proposition}[Path-integrated barrier]
\label{prop:path}
Let $\theta:[0,1]\to\R^p$ be piecewise $C^1$ with $\dot\theta(t)\in\U$, and suppose the metrics along the path satisfy $Q_h^{\theta(t)}(v)\le\bar\rho\,Q_b^{\theta(t)}(v)$ for all $v\in\U$ and $t$. Then
\begin{equation*}
\begin{aligned}
&\bigl\|\Ft_h(\theta(1))-\Ft_h(\theta(0))\bigr\|_2\\
&\qquad\le\ \sqrt{\bar\rho}\int_0^1\Bigl\|\tfrac{d}{dt}\Ft_b(\theta(t))\Bigr\|_2\,dt .
\end{aligned}
\end{equation*}
If instead the damped ratio is bounded by $\bar\rho_\mu$, the right-hand side becomes
\begin{equation*}
\sqrt{\bar\rho_\mu}\Bigl(\int_0^1\bigl\|\tfrac{d}{dt}\Ft_b\bigr\|\,dt+\sqrt\mu\int_0^1\bigl\|\Lambda_{\theta(t)}^{1/2}\dot\theta\bigr\|\,dt\Bigr).
\end{equation*}
\end{proposition}
\begin{proof}
We have
\begin{equation*}
\begin{aligned}
\bigl\|\Ft_h(\theta(1))-\Ft_h(\theta(0))\bigr\|&=\Bigl\|\int_0^1\Ch\dot\theta\,dt\Bigr\|\\
&\le\int_0^1\sqrt{Q_h(\dot\theta)}\,dt\\
&\le\sqrt{\bar\rho}\int_0^1\sqrt{Q_b(\dot\theta)}\,dt,
\end{aligned}
\end{equation*}
and $\sqrt{Q_b(\dot\theta)}=\|\tfrac{d}{dt}\Ft_b(\theta(t))\|$. For the damped case use $\sqrt{x+y}\le\sqrt x+\sqrt y$.
\end{proof}
Any path that moves the safety margins by $\Delta$ must therefore move the benign outputs along a curve of length at least $\Delta/\sqrt{\bar\rho}$, provided the ratio bound holds along the path. This is an assumption, not a consequence of training; \cref{sec:barrier_analysis} tests it empirically. The bound also explains why an amplification term that increases the margin displacement an attack must achieve, such as SEAM's unlearning loss, increases the required benign motion linearly.

% ---------------------------------------------------------
\subsection{Scope}
\label{app:scope}

\Cref{thm:barrier} and \cref{prop:damped,prop:eta,prop:path} concern specified quadratic models or paths along which a ratio bound holds. The trained quantity is the minibatch, damped, normalized $\etak$, differentiated through the Gram with a centered-difference curvature term. None of these results implies robustness to arbitrary harmful fine-tuning, jailbreak resistance, or capability destruction under non-local edits such as weight orthogonalization; those are empirical questions.

% =========================================================
\section{Implementation Details}
\label{app:impl}

\Cref{alg:cbg} lists one defense step; this appendix gives the layer-level details.

\paragraph{Stored tensors.}
The center pass stores, for every linear layer $Z=XW^\top$ in $\U$, the inputs $X$ and the output gradients $\Delta$ of the backward pass of $\mathbf 1^\top\Ft$ (\cref{alg:cbg}). Everything else in the step is formed from these tensors, so the graph is not retained and no second backward pass is run. By direct count they take about 5\,MB per token in bf16 for Llama-2-7B (about 40\,GB at 8{,}192 tokens) and can be offloaded to host memory layer by layer.

\paragraph{Per-example Gram accumulation.}
Let $X_e$ ($T_e\times d_{\mathrm{in}}$) and $\Delta_e$ ($T_e\times d_{\mathrm{out}}$) be the stored tensors on the tokens of example $e$. Because examples do not interact in the forward pass, the block of the per-example gradient $c_e$ belonging to this layer is $\Delta_e^\top X_e$; a safety example sums its two sequences. Inner products in the $\Lam^{-1}$ metric are
\begin{equation*}
K_{ee'}\mathrel{+}=\bigl\langle(\theta\odot\theta+\xi^2)\odot\Delta_e^\top X_e,\;\Delta_{e'}^\top X_{e'}\bigr\rangle ,
\end{equation*}
accumulated in float64. Per-example gradients are formed in chunks of output rows, so only part of one layer is held at a time. Holding all of them for a $14336\times4096$ matrix and $m+n=40$ examples would take about 4.7\,GB in bf16 (9.4\,GB in fp32) by direct count. When $\Lam^{-1}$ is replaced by the identity on a layer, the Gram can be formed without materializing per-example gradients~\citep{li2022large}.

\paragraph{Forming $w$ and $d_\Lambda$.}
After the eigensolve, one more chunked pass over the stored tensors re-forms the per-example gradients and computes, layer by layer, $w=\Lam^{-1}\odot\sum_eb_e\,\Delta_e^\top X_e$ and the diagonal $\sum_ec_e\odot\sum_{e'}G_{ee'}c_{e'}$ that gives $d_\Lambda$ (\cref{prop:gram}). The direction $w$ alone would need only one weighted matrix product per layer; $d_\Lambda$ is what requires the per-example form.

\paragraph{Per-example variant.}
Each linear layer computes $x_e\theta^\top\pm\varepsilon_e\,x_ew_e^\top$ on the tokens of example $e$, with $w_e=\Lam^{-1}C^\top G\mathbf 1_e$ and $\varepsilon_e=r\sqrt{\dim\U/(GKG)_{ee}}$. The $w_e$ of a layer are rebuilt chunk by chunk from the stored tensors in the forward pass and again in the backward pass, where the branch contributes to the input gradient; the scalarization of pass $\sigma$ is $\sum_e\sigma\beta\Ft_e/\varepsilon_e$.

\paragraph{Randomness.}
The sampled targets $\hat y$, dropout masks (disabled by default), sign draws and minibatches are shared between the center pass and both perturbed passes.

\paragraph{Cost (pass-count estimate, not a measurement).}
Let $G_{\mathrm{p}}$ be one forward and backward pass on 4{,}096 tokens. With $m=8$ safety pairs of at most 256 tokens per response and $n=32$ benign sequences of at most 128 tokens, a joint pass costs about $2G_{\mathrm{p}}$. Forming all per-example gradients costs about as much as one weight-gradient computation, roughly a third of a joint pass. One step with the shared direction then costs about $2G_{\mathrm{p}}$ for the center pass, $0.8G_{\mathrm{p}}$ for the Gram, a negligible eigensolve, $0.8G_{\mathrm{p}}$ for $w$ and $d_\Lambda$, $4G_{\mathrm{p}}$ for the two perturbed passes, and $0.3G_{\mathrm{p}}$ for the reference forward of $\Lbgn$: about $8G_{\mathrm{p}}$. SEAM's Algorithm~1 evaluates five to six gradients per step on batches of 8 examples~\citep{wang2026seam}; its token count per batch is not reported, so the ratio depends on it. The per-example variant adds the branch products and the rebuilding of $w_e$ to both perturbed passes, for about $14G_{\mathrm{p}}$ per step. Measured time and memory are reported in \cref{sec:overhead}.

% =========================================================
\section{Hyperparameters}
\label{app:hparams}

\Cref{tab:hparams} lists the CBG settings; values marked TBD are selected on the validation protocol of \cref{sec:setting} and fixed before test attacks. Attack configurations and baseline settings follow SEAM~\citep{wang2026seam}.

\begin{table}[h]
\centering
\small
\caption{CBG hyperparameters.}
\label{tab:hparams}
\begin{tabular}{@{}p{0.40\linewidth}p{0.52\linewidth}@{}}
\toprule
Setting & Value\\
\midrule
Admissible space $\U$ & all linear-layer weights\\
Safety pairs / step $m$ & 8 ($\le256$ tokens per response)\\
Benign sequences / step $n$ & 32 ($\le128$ tokens)\\
Spectral modes $k$ & 4\\
Damping scale $\mu$ & $\kappa_\mu\tr(K_{bb})/n$, differentiated\\
Curvature estimator & shared eigen-sign direction, one draw\\
Loss weights $\alpha,\beta$ & \tbd\\
Damping $\kappa_\mu$, floor $\xi$ & \tbd\\
Perturbation radius $r$ & \tbd\\
Optimizer, lr, schedule, steps & \tbd\\
\bottomrule
\end{tabular}
\end{table}

% =========================================================
\section{Numerical Verification on a Synthetic Model}
\label{app:numerical}

We verified the identities and estimators used above in float64 on two small synthetic models with exact derivatives: a 60-parameter model with an RMS-normalized hidden layer and a normalization gain, which has exact function-preserving rescaling symmetries, and, for the rows marked $\dagger$, a 30-parameter analytic model with closed-form gradients and Hessians ($m=4$, $n=8$, $k=2$). \Cref{tab:numerical} lists the checks. They confirm the algebra; they say nothing about LLM behavior.

\begin{table*}[t]
\centering
\small
\caption{Synthetic checks of the theory (float64).}
\label{tab:numerical}
\begin{tabular}{@{}p{0.22\linewidth}p{0.72\linewidth}@{}}
\toprule
Result & Observation\\
\midrule
\Cref{thm:barrier}(iii) & constrained minimum matches $\gamma^2/(2\Resc)$ to $4\times10^{-12}$ relative (20 random pencils)\\
\Cref{prop:degenerate} & minibatch $\Gbh$ has rank 8 of 60; $\mu\rhoreg$ converges to $\lambda_{\max}(\widetilde C_hP_\perp\widetilde C_h^\top)$ as $\mu\to0$\\
\Cref{prop:damped}(iii) & sandwich holds on 50 random pencils\\
\Cref{prop:reduced} & top-3 eigenvalues and eigenvector agree with the full pencil to machine precision\\
\Cref{prop:reduced}(ii)$^\dagger$ & escape formulas (eigenvector, normalization, both displacements) hold to $1.4\times10^{-15}$\\
\Cref{prop:gram,prop:gradient}$^\dagger$ & Gram-gradient formula matches central differences of $\log\etak$ to $1.6\times10^{-9}$ ($\mu$ differentiated) and $1.5\times10^{-9}$ ($\mu$ held fixed); the alternative forms of $d_\Lambda$ and $\|\Lam^{1/2}w_e\|^2$ agree to $5\times10^{-16}$\\
\Cref{rem:mu}$^\dagger$ & $c_\mu$ matches $\partial\log\etak/\partial\mu$ to $4\times10^{-9}$; fivefold benign inflation leaves $\log\etak$ unchanged ($2\times10^{-16}$) but changes it by $0.41$ with $\mu$ held fixed\\
\Cref{prop:randomized}$^\dagger$ & expectation over all $2^{12}$ sign draws matches the curvature term to $3\times10^{-15}$; a single draw has $200\%$ mean relative error with the eigen-sign pair and $279\%$ with a plain Rademacher pair; the shared and per-example differences are second order (error ratio $4.00$ per halving)\\
\Cref{prop:finite} & gradient error is second order in $\varepsilon$\\
\Cref{prop:reparam} & a function-preserving rescaling changes Euclidean-damped $\rhoreg$ by $+44\%$; $\Lam$-damped $\rhoreg$ and $\eta_1$ unchanged\\
\Cref{prop:jensen} & minibatch mean exceeds the population value by $\approx18\times$, $5\times$, $2\times$ at $m=n=4,16,64$\\
\Cref{prop:mcfisher} & exact expectation equals the token-summed Fisher; KL ratio $0.99997$\\
\Cref{prop:path} & bound holds on a curved path\\
Precision & with bf16-rounded weights the centered difference has $190\%$, $35\%$, $6\%$ error at relative radii $10^{-4},10^{-3},10^{-2}$ (under $0.3\%$ unrounded)\\
\bottomrule
\end{tabular}
\end{table*}

\end{document}